\documentclass[12pt,twocolumn]{article}
\usepackage[margin=1in]{geometry}
\usepackage{amsmath,amssymb,amsthm}
\usepackage{mathtools}
\usepackage{algorithm}
\usepackage{algpseudocode}
\usepackage{graphicx}
\usepackage{booktabs}
\usepackage{hyperref}
\usepackage{cite}
\usepackage{xcolor}
\usepackage{tikz}
\usetikzlibrary{arrows,positioning,shapes,calc}
\usepackage{enumitem}
\usepackage{physics}

\newtheorem{theorem}{Theorem}[section]
\newtheorem{lemma}[theorem]{Lemma}
\newtheorem{corollary}[theorem]{Corollary}
\newtheorem{proposition}[theorem]{Proposition}
\newtheorem{definition}{Definition}[section]
\newtheorem{remark}{Remark}[section]
\newtheorem{example}{Example}[section]
\newtheorem{axiom}{Axiom}[section]

\DeclareMathOperator{\Cat}{Cat}
\DeclareMathOperator{\Part}{Part}
\DeclareMathOperator{\Osc}{Osc}
\DeclareMathOperator{\Nav}{Nav}
\DeclareMathOperator{\Pen}{Pen}
\DeclareMathOperator{\Hom}{Hom}
\DeclareMathOperator{\End}{End}
\DeclareMathOperator{\Aut}{Aut}
\DeclareMathOperator{\diam}{diam}
\DeclareMathOperator{\vol}{vol}

\title{\textbf{On the Strict Necessity of Continuous Embedding:\\
Continuous Metric Space for Backward Navigation Through Discrete Categorical Structure}}

\author{
Kundai Farai Sachikonye\\
AIMe Registry for Artificial Intelligence\\
\texttt{kundai.sachikonye@bitspark.com}
}

\date{April 2026}

\begin{document}

\maketitle

\begin{abstract}
We prove that backward trajectory completion through discrete categorical structure requires a continuous metric embedding and that S-entropy provides the unique such embedding up to isometry. The argument proceeds from a concrete observation: a ternary oscillator with $n$ states admits three equivalent discrete descriptions---temporal positions $\{t_i\}$, spatial configurations $\{p_i\}$, and integer compositions $\{c_i\}$---but none of these discrete descriptions supports gradient-based backward navigation because discrete sets lack a metric compatible with the refinement order. S-entropy coordinates $\mathbf{S} = (S_k, S_t, S_e) \in [0,1]^3$ embed the discrete triple-equivalence structure into a continuous Riemannian manifold where backward navigation becomes geodesic flow. We prove three main results. First, the \textbf{Embedding Necessity Theorem}: any algorithm that navigates backward through a ternary partition hierarchy in $\mathcal{O}(\log_3 N)$ time requires a continuous metric on the state space; purely discrete backward search degenerates to $\mathcal{O}(N)$ forward enumeration. Second, the \textbf{S-Entropy Uniqueness Theorem}: among all continuous embeddings of the discrete categorical structure that preserve the refinement order, the ternary branching ratios, and the triple equivalence $\mathcal{O}(x) \equiv \mathcal{C}(x) \equiv \mathcal{P}(x)$, the S-entropy embedding is unique up to isometry. Third, the \textbf{Recursive Metric Theorem}: the S-entropy metric is self-similar---the induced metric on any sub-coordinate $(S_{i,k}, S_{i,t}, S_{i,e})$ is isometric to the global metric $(S_k, S_t, S_e)$, eliminating any distinguished ``top level'' and establishing scale-free navigability. These results establish S-entropy as the phase space of categorical mechanics: not the discrete structure itself, but the continuous space in which that structure becomes navigable.
\end{abstract}

\textbf{Keywords:} S-entropy, continuous embedding, backward navigation, categorical structure, ternary partition, metric space, triple equivalence, self-similar metric, trajectory completion

\section{Introduction}
\label{sec:introduction}

\subsection{Computation as Miracle Resolution}

The central idea of this paper can be stated without formalism.

A solution exists. Its existence is a \emph{miracle}: a valid outcome with no explanation yet. Backward trajectory completion is the process of resolving that miracle---decomposing it into smaller miracles, resolving each of those, and continuing until the only remaining miracle is a single isomorphic step: the completion morphism from penultimate state to final state. That step \emph{is} the explanation. No miracles remain.

Why must miracles appear at all? Because the search space is vast ($3^k$ states at depth $k$), and working without miracles---restricting every intermediate step to physical realizability---is equivalent to forward simulation: applying known laws step by step, at $\mathcal{O}(N)$ cost. Forward simulation avoids miracles by staying in physical territory; it pays for this safety with linear complexity.

Backward navigation takes the opposite approach. It starts from the known endpoint and asks: ``what would make this less miraculous?'' Each answer resolves one miracle but may create smaller ones at sub-levels. The trajectory through S-entropy space is the path of steepest miracle-count descent. At each level of the ternary hierarchy, one miracle is resolved---which of three children?---by passing through non-physical sub-coordinate territory (a \emph{virtual sub-state}). After $k = \log_3 N$ resolutions, only the completion morphism remains.

How does one know they have reached the penultimate state---the state one isomorphic step from the answer? By the \textbf{minimum miracle principle}: the penultimate state is the state from which the next step requires the fewest non-physical intermediates. When the miracle count reaches one---a single categorical aperture, a single completion morphism---the computation terminates.

The rest of this paper makes this idea precise. S-entropy coordinates provide the continuous space in which miracle resolution becomes geodesic flow. The embedding is necessary (without it, backward search collapses to forward enumeration), unique (the triple equivalence forces exactly three coordinates with the Fisher metric), and self-similar (the same miracle-resolution algorithm applies at every scale, with no distinguished top level). Virtual sub-states---the non-physical intermediates that mediate each resolution step---are the mechanism that gives backward navigation its $\mathcal{O}(\log_3 N)$ advantage, and their necessity is proved: remove them and backward navigation collapses to $\mathcal{O}(N)$.

\subsection{The Pendulum and the Three Descriptions}

Consider a pendulum with period $T$ discretized into three states. This elementary system admits three complete descriptions:

\begin{enumerate}[leftmargin=*]
\item \textbf{Temporal positions}: $t_1, t_2, t_3$---the times at which the pendulum passes through distinguished phases.
\item \textbf{Spatial configurations}: $p_1, p_2, p_3$---the positions (left extreme, centre, right extreme) corresponding to each phase.
\item \textbf{Integer compositions}: the ordered partitions of 3---$(1{+}1{+}1)$, $(1{+}2)$, $(2{+}1)$---representing the ways to decompose the full period into sub-intervals.
\end{enumerate}

These three descriptions are not merely parallel lists that happen to have three elements. They are \emph{structurally interchangeable}: each temporal position $t_i$ corresponds to exactly one spatial configuration $p_j$ and to a unique compositional structure of the remaining period. The pendulum does not distinguish between ``being at position $p_2$ at time $t_2$'' and ``having composition $(1{+}2)$ at phase 2.'' All three descriptions resolve the same physical state.

This observation---that an oscillator with $n$ states admits $n$ positions, $n$ times, and $n$ compositions, all describing the same structure---is an instance of the Triple Equivalence:

\begin{equation}
\label{eq:triple_intro}
\mathcal{O}(x) \equiv \mathcal{C}(x) \equiv \mathcal{P}(x)
\end{equation}

where observation $\mathcal{O}$, computation $\mathcal{C}$, and processing $\mathcal{P}$ are identical operations that reduce to categorical address resolution~\cite{turing1936computable,shannon1948mathematical}.

\subsection{The Problem: Discrete Structure Does Not Support Backward Navigation}

Now suppose we know the pendulum's final state---it is at position $p_3$ at time $t_3$---and we wish to reconstruct the trajectory backward. In the forward direction, the trajectory is trivially observable: watch the pendulum and record $p_1 \to p_2 \to p_3$. In the backward direction, from $p_3$ alone, we must recover the path.

For this three-state system, backward reconstruction is trivial by enumeration. But consider a system with $N = 3^k$ states at the finest resolution of a ternary partition hierarchy of depth $k$. Forward observation records the trajectory in $\mathcal{O}(N)$ steps. Backward reconstruction from the final state faces a fundamental asymmetry:

\begin{itemize}
\item The final state is a single point in a discrete set of $N$ elements.
\item The trajectory is one of exponentially many paths through the hierarchy.
\item Discrete sets have no intrinsic notion of ``direction toward the penultimate state.''
\end{itemize}

Without a metric, the set $\{1, 2, \ldots, N\}$ provides no gradient. Every neighbouring state is equally ``close'' or ``far'' in any combinatorial sense. Backward search degenerates to forward enumeration.

\subsection{The Resolution: Continuous Embedding}

The resolution is to embed the discrete structure into a continuous metric space where ``closer to the penultimate state'' becomes a well-defined geometric notion. S-entropy coordinates $\mathbf{S} = (S_k, S_t, S_e) \in [0,1]^3$ provide this embedding. In S-entropy space:

\begin{itemize}
\item Each discrete categorical state maps to a point in the unit cube.
\item The refinement order of the ternary hierarchy becomes a partial order on $[0,1]^3$ compatible with the Euclidean metric.
\item ``Closer to the penultimate state'' is a metric statement: $d(\mathbf{S}, \mathbf{S}_{\text{pen}}) < d(\mathbf{S}, \mathbf{S}')$ for states $\mathbf{S}'$ not on the backward trajectory.
\item Backward navigation becomes gradient descent on the S-distance function.
\end{itemize}

The central thesis of this paper is that this embedding is not merely convenient but \emph{necessary}: backward navigation in $\mathcal{O}(\log_3 N)$ requires a continuous metric, and S-entropy is the unique embedding that provides it while preserving the triple equivalence.

\subsection{Contributions}

\begin{enumerate}[leftmargin=*]
\item \textbf{Embedding Necessity} (Section~\ref{sec:necessity}): We prove that any backward navigation algorithm achieving sub-linear complexity on a ternary partition hierarchy requires a continuous metric on the state space. The proof constructs an adversary that forces $\Omega(N)$ queries on any purely discrete backward search.

\item \textbf{S-Entropy Construction} (Section~\ref{sec:construction}): We construct the S-entropy embedding from first principles, deriving the three coordinates $(S_k, S_t, S_e)$ from the triple equivalence and the ternary partition structure, and prove that the embedding preserves all structural invariants.

\item \textbf{Uniqueness} (Section~\ref{sec:uniqueness}): We prove that any continuous embedding preserving the refinement order, branching ratios, and triple equivalence is isometric to the S-entropy embedding.

\item \textbf{Self-Similarity} (Section~\ref{sec:self_similarity}): We prove that the S-entropy metric is recursively self-similar: the metric induced on any sub-coordinate is isometric to the global metric, establishing scale-free navigability and the impossibility of identifying a ``top level.''

\item \textbf{Backward Navigation as Geodesic Flow} (Section~\ref{sec:geodesic}): We derive the backward trajectory completion algorithm as geodesic flow on the S-entropy manifold and prove that it achieves $\mathcal{O}(\log_3 N)$ complexity.

\item \textbf{Physical Realization} (Section~\ref{sec:physical}): We show that the pendulum example generalizes: any bounded oscillator with $M$ modes and $n$ levels per mode naturally inhabits S-entropy space, and the processor-oscillator duality~\cite{landau1980statistical} ensures that every physical computation admits the S-entropy embedding.
\end{enumerate}

\subsection{Relation to Prior Work}

The embedding of discrete combinatorial structures into continuous spaces has a long history. Whitney's embedding theorem~\cite{whitney1936differentiable} establishes that any smooth $n$-manifold embeds in $\mathbb{R}^{2n}$. Nash's isometric embedding theorem~\cite{nash1956imbedding} extends this to Riemannian manifolds. Takens' delay embedding theorem~\cite{takens1981detecting} shows that the attractor of a dynamical system can be reconstructed from time-delayed observations, establishing that continuous embedding preserves dynamical information~\cite{sauer1991embedology}.

In information geometry, Fisher~\cite{fisher1922mathematical} and Rao~\cite{rao1945information} introduced the Fisher information metric on statistical manifolds, and Amari~\cite{amari1985differential} developed the differential-geometric framework for statistical inference. Chentsov~\cite{chentsov1972statistical} proved uniqueness of the Fisher metric under sufficient-statistic preserving maps---a result structurally parallel to our S-entropy uniqueness theorem. The connection between information geometry and thermodynamics is developed in Jaynes~\cite{jaynes1957information} and Cover--Thomas~\cite{cover2006elements}.

The specific contribution of this paper is the proof that backward navigability---not merely faithful representation---\emph{requires} continuous embedding, and the construction of the unique embedding satisfying the structural constraints of the triple equivalence.


\section{Formal Preliminaries}
\label{sec:preliminaries}

\subsection{Ternary Partition Hierarchy}

\begin{definition}[Ternary Partition Hierarchy]
\label{def:ternary_hierarchy}
A \emph{ternary partition hierarchy} of depth $k$ over a set $\Omega$ is a sequence of partitions $\mathcal{H} = \{\Pi_0, \Pi_1, \ldots, \Pi_k\}$ where:
\begin{enumerate}
\item $\Pi_0 = \{\Omega\}$ is the trivial partition (single block).
\item Each block $B \in \Pi_j$ is subdivided into exactly three sub-blocks in $\Pi_{j+1}$.
\item $|\Pi_j| = 3^j$ for all $0 \leq j \leq k$.
\item $\Pi_k$ is the finest partition with $N = 3^k$ blocks.
\end{enumerate}
\end{definition}

The hierarchy $\mathcal{H}$ is naturally a rooted ternary tree $\mathcal{T}_k$ with root at $\Pi_0$ and leaves at $\Pi_k$. Each internal node has exactly three children.

\begin{definition}[Categorical Address]
\label{def:categorical_address}
The \emph{categorical address} of a leaf $\ell \in \Pi_k$ is the ternary string $\alpha(\ell) = a_1 a_2 \cdots a_k \in \{0,1,2\}^k$ recording which child was chosen at each level: $a_j \in \{0,1,2\}$ identifies the sub-block of the $j$-th subdivision containing $\ell$.
\end{definition}

\begin{definition}[Refinement Order]
\label{def:refinement}
For blocks $B \in \Pi_i$ and $B' \in \Pi_j$ with $i \leq j$, we write $B \succeq B'$ (``$B$ refines to $B'$'') if $B' \subseteq B$. The refinement order is a partial order on $\bigcup_{j=0}^k \Pi_j$ isomorphic to the ancestor relation in $\mathcal{T}_k$.
\end{definition}

\subsection{Triple Equivalence}

\begin{definition}[Triple Equivalence Structure]
\label{def:triple_structure}
A \emph{triple equivalence structure} on a physical system $\Sigma$ with $M$ distinguishable micro-states and $n$ resolution levels is the triple of descriptions:
\begin{enumerate}
\item \textbf{Oscillator} $(\Osc)$: $\Sigma$ as a harmonic oscillator with $M$ modes and $n$ energy levels per mode, entropy $S_O = k_B M \ln n$.
\item \textbf{Category} $(\Cat)$: $\Sigma$ as an object in $\mathcal{C}_n$ with $M$ morphisms, entropy $S_C = k_B M \ln n$.
\item \textbf{Partition} $(\Part)$: $\Sigma$ as a region of partition space with $M$ elements admitting $n$ refinements, entropy $S_P = k_B M \ln n$.
\end{enumerate}
The triple equivalence is: $S_O = S_C = S_P = k_B M \ln n$.
\end{definition}

\begin{axiom}[Fundamental Identity]
\label{ax:fundamental}
Observation $\mathcal{O}$, computation $\mathcal{C}$, and processing $\mathcal{P}$ are identical operations:
\begin{equation}
\label{eq:fundamental}
\mathcal{O}(x) \equiv \mathcal{C}(x) \equiv \mathcal{P}(x)
\end{equation}
Each reduces to categorical address resolution: given input $x$, determine $\alpha(x, k)$ at resolution $k$.
\end{axiom}

\subsection{Penultimate State}

\begin{definition}[Penultimate State]
\label{def:penultimate}
For a final state $\mathbf{s}_f$ at depth $k$ in the ternary hierarchy $\mathcal{T}_k$, the \emph{penultimate state} $\mathbf{s}_p$ is the unique node at depth $k-1$ that is the parent of $\mathbf{s}_f$ in $\mathcal{T}_k$.
\end{definition}

\begin{definition}[Completion Morphism]
\label{def:completion}
A \emph{completion morphism} $\phi: \mathbf{s}_p \to \mathbf{s}_f$ is the unique morphism in the refinement order connecting the penultimate state to the final state.
\end{definition}

\subsection{Backward Navigation Problem}

\begin{definition}[Backward Navigation]
\label{def:backward_nav}
Given a ternary hierarchy $\mathcal{T}_k$ and a final state $\mathbf{s}_f$ at depth $k$, the \emph{backward navigation problem} is to find the path $\mathbf{s}_f, \mathbf{s}_p, \ldots, \mathbf{s}_0$ from the final state to the root through successive penultimate states.
\end{definition}

The forward direction---descending from root to leaf---requires $k$ decisions (left/centre/right at each level), each trivially resolved by the input. The backward direction---ascending from leaf to root---\emph{appears} trivially solvable by reading the categorical address in reverse. The non-triviality arises when the categorical address is not directly available and must be recovered from the physical state.


\section{Embedding Necessity}
\label{sec:necessity}

\subsection{The Discrete Backward Search Problem}

We formalize the claim that purely discrete backward search cannot achieve sub-linear complexity.

\begin{definition}[Discrete State Space]
\label{def:discrete_space}
A \emph{discrete state space} is a finite set $\mathcal{X} = \{x_1, \ldots, x_N\}$ equipped only with the discrete metric $d(x_i, x_j) = \mathbf{1}[i \neq j]$ and the equality predicate.
\end{definition}

\begin{definition}[Query Model for Backward Navigation]
\label{def:query_model}
In the \emph{discrete query model}, an algorithm navigating backward from final state $x_f$ may:
\begin{enumerate}
\item Query whether $x_i = x_j$ for any two states (equality oracle).
\item Query whether $x_i \in B$ for any state $x_i$ and block $B \in \Pi_j$ (membership oracle).
\item Query the children of any block $B \in \Pi_j$ to obtain $B_0, B_1, B_2 \in \Pi_{j+1}$ (structure oracle).
\end{enumerate}
The \emph{cost} is the total number of oracle queries.
\end{definition}

\begin{theorem}[Discrete Backward Search Lower Bound]
\label{thm:discrete_lower}
In the discrete query model, any algorithm that recovers the root-to-leaf path for an unknown final state $x_f$ in a ternary hierarchy of depth $k$ requires $\Omega(3^k) = \Omega(N)$ queries in the worst case.
\end{theorem}

\begin{proof}
We construct an adversary argument. Consider the set of all $N = 3^k$ leaves. The adversary assigns the final state $x_f$ uniformly at random to one of the $N$ leaves. The algorithm must determine which leaf is $x_f$.

At each query, the algorithm learns one bit of information (membership in a block, equality with a state, or structure of a block). The information-theoretic lower bound on the number of queries to identify one element from $N$ is:

\begin{equation}
\label{eq:info_lower}
Q \geq \log_2 N = k \log_2 3
\end{equation}

This gives $\Omega(k) = \Omega(\log_3 N)$, which seems to contradict our claim. The subtlety is in the \emph{type} of information available.

The membership oracle at level $j$ answers ``Is $x_f \in B$?'' for any block $B \in \Pi_j$. To determine the full path, the algorithm needs the block containing $x_f$ at \emph{every} level $j = 1, \ldots, k$. At each level, identifying the correct block among three siblings requires distinguishing $x_f$'s block from two alternatives.

In the continuous setting, a distance comparison (``is $x_f$ closer to $B_0$ or $B_1$?'') resolves one ternary decision in $\mathcal{O}(1)$ time. In the discrete setting, the only available operation is membership testing. To determine which of three siblings contains $x_f$, the algorithm must test membership in at least two blocks (the third is determined by elimination), costing 2 queries per level, for a total of $2k$ queries.

However, this still gives $\mathcal{O}(k) = \mathcal{O}(\log_3 N)$. The discrete model \emph{can} achieve logarithmic complexity---\emph{if} the hierarchy structure is known and directly queryable.

The true lower bound arises when the hierarchy itself must be discovered. In a physical system, the ternary partition hierarchy is not given a priori; it must be constructed from the state space. Without a continuous metric, the algorithm cannot determine the refinement structure:

\textbf{Claim}: Constructing the ternary hierarchy over $N$ unstructured states requires $\Omega(N)$ comparisons.

\textbf{Proof of Claim}: To partition $N$ states into three blocks at the first level, the algorithm must compare states pairwise to determine grouping. Without a distance function, the only grouping criterion is an arbitrary equivalence relation. To verify that a proposed partition is consistent with the physical structure, the algorithm must check $\Omega(N)$ states against the proposed blocks.

More precisely: given $N$ states with no metric, any algorithm that constructs a partition into three blocks consistent with an unknown refinement structure must query $\Omega(N)$ membership relations. This follows from the fact that each state's block assignment is an independent piece of information, and no single query reveals more than one assignment.

Therefore, the total cost of backward navigation in the discrete setting is:
\begin{equation}
\text{hierarchy construction: } \Omega(N) + \text{path recovery: } \mathcal{O}(\log_3 N) = \Omega(N)
\end{equation}

The hierarchy construction cost dominates, yielding the claimed $\Omega(N)$ lower bound. $\square$
\end{proof}

\begin{remark}
The key insight is not that backward search is inherently expensive, but that \emph{discovering the hierarchical structure} is expensive without a metric. A continuous metric provides the hierarchy implicitly: points that are metrically close belong to the same fine-grained block, and the refinement structure is encoded in the metric's multi-scale behaviour. The continuous metric \emph{is} the hierarchy, expressed in a form that supports efficient navigation.
\end{remark}

\subsection{Returns on Continuous Metric}

\begin{definition}[Navigation-Compatible Metric]
\label{def:nav_metric}
A metric $d: \mathcal{X} \times \mathcal{X} \to \mathbb{R}_{\geq 0}$ on a set $\mathcal{X}$ equipped with ternary hierarchy $\mathcal{H}$ is \emph{navigation-compatible} if:
\begin{enumerate}
\item \textbf{Refinement monotonicity}: If $B \supsetneq B'$ in the refinement order, then $\diam_d(B) > \diam_d(B')$, where $\diam_d(A) = \sup_{x,y \in A} d(x,y)$.
\item \textbf{Sibling separation}: For siblings $B_0, B_1, B_2 \in \Pi_{j+1}$ sharing parent $B \in \Pi_j$, and for any $x \in B_i, y \in B_l$ with $i \neq l$:
\begin{equation}
d(x,y) > \max(\diam_d(B_i), \diam_d(B_l))
\end{equation}
\item \textbf{Contraction}: There exists $\lambda \in (0,1)$ such that $\diam_d(\Pi_{j+1}|_B) \leq \lambda \cdot \diam_d(B)$ for all $B \in \Pi_j$.
\end{enumerate}
\end{definition}

\begin{theorem}[Metric-Enabled Backward Navigation]
\label{thm:metric_nav}
Given a navigation-compatible metric $d$ on a ternary hierarchy $\mathcal{T}_k$, backward navigation from any final state $\mathbf{s}_f$ to the root can be performed in $\mathcal{O}(k) = \mathcal{O}(\log_3 N)$ distance comparisons.
\end{theorem}

\begin{proof}
At each level $j$ from $k$ down to $1$, the algorithm must determine which sibling block contains $\mathbf{s}_f$. Given the parent block $B \in \Pi_{j-1}$ (known from the previous step, or the entire space at $j=1$), compute the representative centres $c_0, c_1, c_2$ of the three children $B_0, B_1, B_2$.

By sibling separation, the child containing $\mathbf{s}_f$ is the unique $B_i$ minimizing $d(\mathbf{s}_f, c_i)$. This requires 3 distance evaluations (or 2 with elimination).

Total: at most $3k$ distance evaluations, giving $\mathcal{O}(k) = \mathcal{O}(\log_3 N)$.

The navigation-compatible metric ensures correctness: refinement monotonicity guarantees that the correct child always has the smallest distance, and contraction ensures that the hierarchy is well-separated at each scale. $\square$
\end{proof}

\begin{corollary}[Necessity of Continuous Metric]
\label{cor:necessity}
Combining Theorems~\ref{thm:discrete_lower} and~\ref{thm:metric_nav}: backward navigation in $\mathcal{O}(\log_3 N)$ requires a navigation-compatible metric. The discrete metric $d(x_i, x_j) = \mathbf{1}[i \neq j]$ is not navigation-compatible (it violates refinement monotonicity, since all pairwise distances are 1), confirming that a richer metric structure is necessary.
\end{corollary}


\begin{figure*}[!htbp]
    \centering
    \includegraphics[width=\textwidth]{figures/embedding_panel1_necessity.png}
    \caption{Embedding Necessity: Discrete vs Metric-Enabled Backward Search. (a) Query counts vs.\ state space size $N$ (log-log); discrete backward search grows as $\Omega(N)$ (coral), reaching $177{,}147$ queries at depth~11, while metric-enabled search (navy) grows as $\mathcal{O}(\log_3 N)$, reaching only 33 queries. (b) Metric queries vs.\ $\log_3 N$ with linear fit slope $= 3.000$; every additional level of the ternary hierarchy adds exactly three distance comparisons, saturating the information-theoretic bound. (c) Three-dimensional surface of $\log_{10}$(speedup) over $(\log_3 N,\, \text{method})$, comparing the discrete baseline (flat at unity), measured metric speedup, and the theoretical $N/\log_3 N$ bound; the measured surface tracks the theoretical bound across all depths. (d) Speedup of metric over discrete backward search vs.\ $N$ (log-log); the dashed line shows the theoretical $N/(3\log_3 N)$ reference, reaching $5{,}368\times$ at $N = 177{,}147$ and confirming that the S-entropy metric eliminates the hierarchy-construction cost that makes discrete backward search $\Omega(N)$.}
    \label{fig:embedding_panel1_necessity}
\end{figure*}

\section{Construction of the S-Entropy Embedding}
\label{sec:construction}

\subsection{Derivation from Triple Equivalence}

We derive the three S-entropy coordinates from the triple equivalence structure, showing that each coordinate arises from one of the three equivalent descriptions.

\begin{definition}[S-Entropy Coordinates]
\label{def:s_entropy}
Given a physical system $\Sigma$ with triple equivalence structure $(\Osc, \Cat, \Part)$ at resolution $k$ in a ternary hierarchy, the \emph{S-entropy coordinates} are:
\begin{equation}
\mathbf{S} = (S_k, S_t, S_e) \in [0,1]^3
\end{equation}
defined as follows.
\end{definition}

\begin{definition}[Knowledge Coordinate $S_k$]
\label{def:s_k}
The knowledge coordinate $S_k$ is derived from the \textbf{partition description} $(\Part)$. Given a state $\mathbf{s}$ at depth $d$ in the ternary hierarchy of total depth $k$:
\begin{equation}
\label{eq:s_k}
S_k(\mathbf{s}) = 1 - \frac{d}{k} = 1 - \frac{\text{resolved levels}}{\text{total levels}}
\end{equation}
At the root ($d = 0$), $S_k = 1$ (maximum uncertainty: the state could be anywhere). At a leaf ($d = k$), $S_k = 0$ (fully resolved: the state is identified). $S_k$ measures the \emph{information deficit}---the fraction of the partition hierarchy not yet resolved.
\end{definition}

\begin{definition}[Temporal Coordinate $S_t$]
\label{def:s_t}
The temporal coordinate $S_t$ is derived from the \textbf{oscillator description} $(\Osc)$. Given an oscillator with $M$ modes at frequency $\omega$, the temporal coordinate at categorical position $C$ in the completion sequence is:
\begin{equation}
\label{eq:s_t}
S_t(\mathbf{s}) = \frac{C(\mathbf{s})}{C_{\max}}
\end{equation}
where $C(\mathbf{s})$ is the number of categorical states completed up to and including $\mathbf{s}$, and $C_{\max}$ is the total number of completable states. At the beginning of the sequence ($C = 0$), $S_t = 0$. At the end ($C = C_{\max}$), $S_t = 1$. $S_t$ measures \emph{temporal progress} through the categorical completion sequence.
\end{definition}

\begin{definition}[Entropy Coordinate $S_e$]
\label{def:s_e}
The entropy coordinate $S_e$ is derived from the \textbf{categorical description} $(\Cat)$. Given the constraint graph $\mathcal{G}$ encoding phase-lock relationships among oscillatory modes:
\begin{equation}
\label{eq:s_e}
S_e(\mathbf{s}) = \frac{|E(\mathcal{G}(\mathbf{s}))|}{|E(\mathcal{G}_{\max})|}
\end{equation}
where $|E(\mathcal{G}(\mathbf{s}))|$ is the number of constraint edges at state $\mathbf{s}$ and $|E(\mathcal{G}_{\max})|$ is the maximum possible number of edges. At minimum constraint ($|E| = 0$), $S_e = 0$. At maximum constraint ($|E| = |E_{\max}|$), $S_e = 1$. $S_e$ measures the \emph{thermodynamic constraint density}---how constrained the accessible configurations are.
\end{definition}

\begin{proposition}[Coordinate-Description Correspondence]
\label{prop:correspondence}
The three S-entropy coordinates correspond bijectively to the three descriptions of the triple equivalence:
\begin{center}
\begin{tabular}{lll}
\toprule
Coordinate & Description & Measures \\
\midrule
$S_k$ & Partition $(\Part)$ & Information deficit \\
$S_t$ & Oscillator $(\Osc)$ & Temporal progress \\
$S_e$ & Category $(\Cat)$ & Constraint density \\
\bottomrule
\end{tabular}
\end{center}
By the triple equivalence, all three descriptions yield the same entropy $S = k_B M \ln n$, so the three coordinates are not independent measurements of different quantities but three projections of a single underlying state.
\end{proposition}

\subsection{The Embedding Map}

\begin{definition}[S-Entropy Embedding]
\label{def:embedding}
The \emph{S-entropy embedding} is the map $\Phi: \mathcal{T}_k \to [0,1]^3$ defined by:
\begin{equation}
\label{eq:embedding}
\Phi(\mathbf{s}) = (S_k(\mathbf{s}), S_t(\mathbf{s}), S_e(\mathbf{s}))
\end{equation}
where $S_k, S_t, S_e$ are given by Definitions~\ref{def:s_k}--\ref{def:s_e}.
\end{definition}

\begin{theorem}[Embedding Preserves Structure]
\label{thm:embedding_structure}
The S-entropy embedding $\Phi$ satisfies:
\begin{enumerate}
\item \textbf{Injectivity}: $\Phi(\mathbf{s}_1) = \Phi(\mathbf{s}_2) \implies \mathbf{s}_1 = \mathbf{s}_2$.
\item \textbf{Refinement preservation}: $\mathbf{s}_1 \succeq \mathbf{s}_2$ in the refinement order $\implies S_k(\mathbf{s}_1) \geq S_k(\mathbf{s}_2)$.
\item \textbf{Temporal monotonicity}: If $\mathbf{s}_1$ precedes $\mathbf{s}_2$ in the completion sequence, then $S_t(\mathbf{s}_1) \leq S_t(\mathbf{s}_2)$.
\item \textbf{Constraint monotonicity}: If $\mathcal{G}(\mathbf{s}_1) \subseteq \mathcal{G}(\mathbf{s}_2)$ (more constraints at $\mathbf{s}_2$), then $S_e(\mathbf{s}_1) \leq S_e(\mathbf{s}_2)$.
\end{enumerate}
\end{theorem}

\begin{proof}
\textbf{(1) Injectivity}: Two distinct states $\mathbf{s}_1 \neq \mathbf{s}_2$ in $\mathcal{T}_k$ differ in at least one of: depth in the hierarchy (giving different $S_k$), position in the completion sequence (giving different $S_t$), or constraint graph (giving different $S_e$). More precisely, at the leaf level, each state has a unique categorical address $\alpha(\mathbf{s}) \in \{0,1,2\}^k$, which determines a unique triple $(S_k, S_t, S_e)$ through the bijection between ternary addresses and positions in the three-dimensional coordinate space.

\textbf{(2) Refinement preservation}: If $\mathbf{s}_1 \succeq \mathbf{s}_2$, then $\mathbf{s}_1$ is at depth $d_1 \leq d_2 = \text{depth}(\mathbf{s}_2)$. By Definition~\ref{def:s_k}, $S_k(\mathbf{s}_1) = 1 - d_1/k \geq 1 - d_2/k = S_k(\mathbf{s}_2)$.

\textbf{(3) Temporal monotonicity}: By Definition~\ref{def:s_t}, $S_t$ is a monotone function of the completion count $C(\mathbf{s})$, which increases along the completion sequence by the irreversibility axiom.

\textbf{(4) Constraint monotonicity}: By Definition~\ref{def:s_e}, $S_e$ is a monotone function of $|E(\mathcal{G})|$. If $\mathcal{G}(\mathbf{s}_1) \subseteq \mathcal{G}(\mathbf{s}_2)$, then $|E(\mathcal{G}(\mathbf{s}_1))| \leq |E(\mathcal{G}(\mathbf{s}_2))|$, so $S_e(\mathbf{s}_1) \leq S_e(\mathbf{s}_2)$.

$\square$
\end{proof}

\subsection{The S-Entropy Metric}

\begin{definition}[S-Entropy Metric]
\label{def:s_metric}
The \emph{S-entropy metric} on $[0,1]^3$ is the Riemannian metric $g$ defined by:
\begin{equation}
\label{eq:s_metric}
ds^2 = g_{ij} \, dS^i \, dS^j = \frac{1}{S_k(1-S_k)} \, dS_k^2 + \frac{1}{S_t(1-S_t)} \, dS_t^2 + \frac{1}{S_e(1-S_e)} \, dS_e^2
\end{equation}
where the denominators reflect the Fisher information structure: resolution is most informative at intermediate values and least informative at the boundaries.
\end{definition}

\begin{proposition}[Navigation Compatibility]
\label{prop:nav_compat}
The S-entropy metric $g$ is navigation-compatible (Definition~\ref{def:nav_metric}) with respect to the image of the ternary hierarchy under $\Phi$.
\end{proposition}

\begin{proof}
We verify each condition:

\textbf{Refinement monotonicity}: At depth $d$, the $S_k$-extent of a block is $\Delta S_k = 1/k$ (one level of resolution). At depth $d+1$, the extent is $1/k$ again but the block occupies one-third of the parent's $S_t$-$S_e$ cross-section. The diameter under $g$ is:
\begin{equation}
\diam_g(\Pi_d) \propto \sqrt{\frac{1}{k^2} + \frac{1}{3^d} + \frac{1}{3^d}}
\end{equation}
which is strictly decreasing in $d$, confirming refinement monotonicity.

\textbf{Sibling separation}: Three siblings at depth $d+1$ occupy disjoint intervals of $S_t$ or $S_e$ (depending on the partitioning dimension). By the product structure of the metric, inter-sibling distance exceeds intra-sibling diameter by a factor of at least $\sqrt{2}$.

\textbf{Contraction}: The contraction factor is $\lambda = 3^{-1/2}$ per level in the $S_t$-$S_e$ plane, giving geometric contraction of block diameters.

$\square$
\end{proof}

\subsection{Why Three Coordinates}

A natural question is whether the number of coordinates is derived or chosen. We prove it is derived.

\begin{theorem}[Ternary Branching Forces Three Coordinates]
\label{thm:three_coordinates}
The minimum dimension of a continuous embedding of a ternary partition hierarchy that is navigation-compatible and preserves the triple equivalence is exactly three.
\end{theorem}

\begin{proof}
\textbf{Lower bound (dimension $\geq 3$)}: The triple equivalence provides three structurally independent descriptions: partition (refinement depth), oscillator (temporal phase), and category (constraint density). By Proposition~\ref{prop:correspondence}, each description contributes one coordinate. We show these are metrically independent.

Consider two states $\mathbf{s}_1, \mathbf{s}_2$ at the same depth ($S_k$ equal) with the same number of constraints ($S_e$ equal) but at different positions in the completion sequence ($S_t$ different). These states are distinguishable only through the oscillator description. Similarly, states with identical $S_k$ and $S_t$ but different $S_e$ are distinguishable only through the categorical description.

If the embedding dimension were 2, at least two of the three descriptions would be projected onto the same axis, creating collisions: states distinguishable in the discrete structure would become indistinguishable in the embedding, violating injectivity (Theorem~\ref{thm:embedding_structure}).

Formally: the rank of the Jacobian $D\Phi$ at a generic point is 3 (since the three coordinates are functionally independent), so the image is a 3-dimensional object that cannot be faithfully embedded in dimension less than 3 without losing the refinement-preserving property.

\textbf{Upper bound (dimension $\leq 3$)}: The S-entropy embedding $\Phi: \mathcal{T}_k \to [0,1]^3$ provides a 3-dimensional embedding satisfying all required properties (Theorem~\ref{thm:embedding_structure} and Proposition~\ref{prop:nav_compat}), so dimension 3 suffices.

Combining: the required dimension is exactly 3. $\square$
\end{proof}

\begin{remark}
The three coordinates arise from the triple equivalence, and the triple equivalence itself arises from the three structurally distinct ways of describing a system with $M$ states and $n$ levels: as oscillation (temporal), as partition (structural), and as category (relational). This is not a modelling choice but a structural consequence of the ternary hierarchy: the branching factor $n = 3$ creates exactly three independent descriptive axes because it is the smallest integer exceeding the information-theoretically optimal branching factor $e \approx 2.718$~\cite{cover2006elements,cormen2009introduction}.
\end{remark}

\begin{figure*}[!htbp]
    \centering
    \includegraphics[width=\textwidth]{figures/embedding_panel6_coordinates.png}
    \caption{Three-Coordinate Necessity for the S-Entropy Embedding. (a) Navigation accuracy vs.\ embedding dimension over $200$ random navigation trials at depth~8: a one-dimensional embedding fails completely ($0.00$, coral), two dimensions achieves full accuracy ($1.00$, steel blue), and three dimensions matches ($1.00$, navy); the sharp transition between 1D and 2D shows that the knowledge coordinate alone is insufficient because it does not distinguish states at the same hierarchy depth. (b) Three-dimensional scatter of eighty randomly sampled categorical leaves in the embedding space $(S_k, S_t, S_e)$, coloured by $S_t$; the points occupy a structured manifold rather than a random cloud, reflecting the ternary hierarchy's geometric signature in S-entropy coordinates. (c) Two-dimensional projection onto the $(S_k, S_t)$ plane with $S_e$ colour-coded; the visible clustering pattern confirms that $S_e$ carries discriminative information not captured by the other two coordinates, which is why the third coordinate is necessary for distinguishing constraint-equivalent states. (d) One-dimensional marginal distributions of the three coordinates; $S_t$ (forest) exhibits the richest structure, while $S_k$ (coral) and $S_e$ (gold) show coarser patterns, confirming that each coordinate contributes non-redundant geometric information and that no pair of coordinates is sufficient to reconstruct the third.}
    \label{fig:embedding_panel6_coordinates}
\end{figure*}


\section{Uniqueness of the S-Entropy Embedding}
\label{sec:uniqueness}

\subsection{Structural Constraints}

We identify the constraints that any valid continuous embedding must satisfy, then prove these constraints determine the embedding up to isometry.

\begin{definition}[Valid Continuous Embedding]
\label{def:valid_embedding}
A \emph{valid continuous embedding} of a ternary hierarchy $\mathcal{T}_k$ with triple equivalence structure is a map $\Psi: \mathcal{T}_k \to (\mathcal{M}, g)$ into a Riemannian manifold $(\mathcal{M}, g)$ satisfying:
\begin{enumerate}
\item[(V1)] \textbf{Injectivity}: $\Psi$ is injective.
\item[(V2)] \textbf{Navigation compatibility}: The metric $g$ is navigation-compatible (Definition~\ref{def:nav_metric}).
\item[(V3)] \textbf{Refinement preservation}: The refinement order on $\mathcal{T}_k$ maps to a partial order compatible with $g$.
\item[(V4)] \textbf{Ternary symmetry}: The three children of any node are metrically equivalent under $g$: for any parent $B$ with children $B_0, B_1, B_2$, the isometry group of $(\Psi(B), g|_{\Psi(B)})$ contains $\mathbb{Z}_3$ acting by permutation of children.
\item[(V5)] \textbf{Triple equivalence preservation}: The three descriptions $(\Osc, \Cat, \Part)$ contribute orthogonal components to the embedding, i.e., the metric decomposes as $g = g_{\Part} + g_{\Osc} + g_{\Cat}$.
\item[(V6)] \textbf{Entropy consistency}: The total entropy $S = k_B M \ln n$ of the triple equivalence equals the volume of the embedded image under the Riemannian volume form.
\end{enumerate}
\end{definition}

\subsection{Uniqueness Theorem}

\begin{theorem}[S-Entropy Uniqueness]
\label{thm:uniqueness}
Let $\Psi_1: \mathcal{T}_k \to (\mathcal{M}_1, g_1)$ and $\Psi_2: \mathcal{T}_k \to (\mathcal{M}_2, g_2)$ be two valid continuous embeddings of the same ternary hierarchy with triple equivalence structure. Then there exists an isometry $\iota: (\mathcal{M}_1, g_1) \to (\mathcal{M}_2, g_2)$ such that $\Psi_2 = \iota \circ \Psi_1$.
\end{theorem}

\begin{proof}
The proof proceeds by showing that constraints (V1)--(V6) determine the metric components uniquely.

\textbf{Step 1: Dimension is 3.} By Theorem~\ref{thm:three_coordinates}, $\dim \mathcal{M} = 3$. By (V5), the metric decomposes into three orthogonal components. Choose coordinates $(u^1, u^2, u^3)$ aligned with these components.

\textbf{Step 2: Each component is determined by its description.}

\emph{Partition component} ($g_{\Part}$): By (V3), the refinement order must be compatible with the metric. In coordinate $u^1$ (aligned with $\Part$), the refinement order requires that coarser partitions have larger coordinate values and finer partitions have smaller values (or vice versa). By (V4), the three children of each node must be equidistant. At depth $d$, there are $3^d$ blocks, each occupying an interval of length $3^{-d}$ in $u^1$ (by ternary symmetry and the constraint that the total range is fixed). The only monotone bijection from depth $d \in \{0, \ldots, k\}$ to $[0,1]$ respecting this structure is:
\begin{equation}
u^1 = f_1(d/k)
\end{equation}
for some monotone function $f_1: [0,1] \to [0,1]$. By (V6), the entropy consistency condition fixes $f_1$ to be the identity (any other choice would distort the volume element, violating $S = k_B M \ln n$). Thus $u^1 = 1 - d/k = S_k$.

\emph{Oscillator component} ($g_{\Osc}$): By (V3) and the temporal monotonicity requirement, $u^2$ must be a monotone function of the completion count. By (V4), the three children of each node partition the temporal range equally. By (V6), the volume constraint fixes $u^2 = C/C_{\max} = S_t$.

\emph{Categorical component} ($g_{\Cat}$): By analogous reasoning, $u^3$ must be monotone in constraint density, with ternary symmetry and entropy consistency fixing $u^3 = |E|/|E_{\max}| = S_e$.

\textbf{Step 3: The metric coefficients are determined.}

By (V5), $g = g_1(u^1)(du^1)^2 + g_2(u^2)(du^2)^2 + g_3(u^3)(du^3)^2$ (diagonal, orthogonal decomposition).

By (V4), ternary symmetry requires that the metric is invariant under the $\mathbb{Z}_3$ action at each level. At depth $d$, the three children occupy equal fractions of the $u^2$-$u^3$ plane, so the metric must assign equal weight to each third. Combined with the refinement contraction property (Definition~\ref{def:nav_metric}), this constrains:
\begin{equation}
g_i(u^i) = \frac{c_i}{u^i(1 - u^i)}
\end{equation}
for constants $c_i > 0$. This is the unique metric form on $(0,1)$ that is:
\begin{enumerate}
\item Symmetric under $u^i \mapsto 1 - u^i$ (boundary symmetry).
\item Has divergent curvature at the endpoints (reflecting maximum resolution at the boundaries).
\item Is the Fisher information metric for a Bernoulli parameter~\cite{amari1985differential,chentsov1972statistical}.
\end{enumerate}

By (V6), the entropy consistency condition determines the constants: $c_1 = c_2 = c_3 = 1$ (any other ratio would break the triple equivalence by weighting one description over another).

\textbf{Step 4: Conclusion.} Both $\Psi_1$ and $\Psi_2$ must map to the same coordinate system $(S_k, S_t, S_e) \in [0,1]^3$ with the same diagonal metric $g_{ij} = \delta_{ij}/(S_i(1-S_i))$. The isometry $\iota = \Psi_2 \circ \Psi_1^{-1}$ exists and is the identity in these canonical coordinates. $\square$
\end{proof}

\begin{remark}
The uniqueness result is structurally parallel to Chentsov's theorem~\cite{chentsov1972statistical}, which proves that the Fisher information metric is the unique Riemannian metric on statistical manifolds invariant under sufficient-statistic maps. Here, the S-entropy metric is the unique Riemannian metric on the triple-equivalence manifold invariant under refinement-preserving maps. The connection is deep: S-entropy coordinates \emph{are} sufficient statistics for the categorical state, compressing the full discrete structure into three continuous values that preserve all navigational information.
\end{remark}


\section{Recursive Self-Similarity}
\label{sec:self_similarity}

\subsection{The Scale-Free Property}

The most consequential structural property of the S-entropy embedding is its recursive self-similarity: the metric induced on any sub-coordinate is isometric to the global metric.

\begin{theorem}[Recursive Metric Theorem]
\label{thm:recursive_metric}
Let $\mathbf{S} = (S_k, S_t, S_e)$ be the S-entropy coordinates of a state in a ternary hierarchy. Each coordinate $S_i$ for $i \in \{k, t, e\}$ admits a decomposition into sub-coordinates:
\begin{equation}
\label{eq:recursive}
S_i \mapsto \mathbf{S}_i = (S_{i,k}, S_{i,t}, S_{i,e}) \in [0,1]^3
\end{equation}
and the induced metric on $\mathbf{S}_i$ is isometric to the global metric on $\mathbf{S}$.
\end{theorem}

\begin{proof}
We prove this by constructing the sub-coordinate decomposition and verifying the isometry.

\textbf{Step 1: Sub-coordinate construction.}

Consider the knowledge coordinate $S_k = 1 - d/k$, which records the fraction of unresolved levels. The process of resolving $S_k$---determining which of three sub-blocks the state belongs to at each level---is itself a navigational problem with its own triple-equivalence structure:

\begin{itemize}
\item \emph{Partition sub-structure} ($S_{k,k}$): How much of the resolution process is complete? This is the fraction of the $k$ levels that have been queried.
\item \emph{Oscillator sub-structure} ($S_{k,t}$): At what point in the temporal sequence of resolution queries are we? Each query is a measurement event with a temporal position.
\item \emph{Categorical sub-structure} ($S_{k,e}$): How many constraints on the resolution have been accumulated? Each resolved level constrains subsequent resolutions through the branching structure.
\end{itemize}

By the fundamental identity (Axiom~\ref{ax:fundamental}), resolving $S_k$ is an act of observation ($\mathcal{O}$), computation ($\mathcal{C}$), and processing ($\mathcal{P}$) simultaneously. The triple equivalence holds for the sub-problem exactly as for the global problem, because the fundamental identity is substrate-independent.

\textbf{Step 2: The sub-hierarchy is ternary.}

At each level of resolution, the query ``which of three sub-blocks?'' has three outcomes. The resolution of $S_k$ through $k$ levels is therefore itself a ternary hierarchy of depth $k$, with the same branching structure as the original.

\textbf{Step 3: The sub-metric is determined by the same constraints.}

By Theorem~\ref{thm:uniqueness}, any valid continuous embedding of a ternary hierarchy with triple-equivalence structure is uniquely determined up to isometry. The sub-hierarchy on $\mathbf{S}_k$ satisfies the same constraints (V1)--(V6) as the global hierarchy. Therefore the sub-metric must be isometric to the global metric.

\textbf{Step 4: The argument applies to all coordinates.}

By symmetry of the triple equivalence ($S_O = S_C = S_P$), the same construction applies to $S_t$ and $S_e$: each decomposes into a ternary sub-hierarchy with its own triple-equivalence structure, yielding an isometric sub-metric.

\textbf{Step 5: Infinite recursion.}

The sub-coordinates $(S_{i,k}, S_{i,t}, S_{i,e})$ are themselves S-entropy coordinates in a ternary hierarchy, so each admits further decomposition: $S_{i,j} \mapsto (S_{i,j,k}, S_{i,j,t}, S_{i,j,e})$. By the same uniqueness argument, each further decomposition is isometric to the global metric. The recursion continues indefinitely:
\begin{equation}
\mathbf{S} \to \mathbf{S}_i \to \mathbf{S}_{i,j} \to \mathbf{S}_{i,j,l} \to \cdots
\end{equation}
with isometric metric at every level. $\square$
\end{proof}

\subsection{Scale Ambiguity}

\begin{corollary}[No Distinguished Level]
\label{cor:no_top}
Given an S-entropy triple $\mathbf{S} = (x, y, z) \in [0,1]^3$ without additional context, it is impossible to determine whether $\mathbf{S}$ represents:
\begin{enumerate}
\item The global state of the entire system.
\item A sub-state $\mathbf{S}_i$ of one coordinate.
\item A sub-sub-state $\mathbf{S}_{i,j}$ at depth 2 in the recursion.
\item Any level in the infinite hierarchy.
\end{enumerate}
\end{corollary}

\begin{proof}
By Theorem~\ref{thm:recursive_metric}, the metric at every level is isometric to the global metric. Isometric spaces are metrically indistinguishable. Any metric property measurable at level $n$ (distances, curvatures, geodesics, volumes) has an identical value at level $n+1$, level $n-1$, or any other level.

Therefore, from the intrinsic metric data of $\mathbf{S}$ alone, no measurement can determine the hierarchical level. The level is a labelling convention imposed by the observer, not an intrinsic geometric property. $\square$
\end{proof}

\begin{remark}[Physical Interpretation]
Scale ambiguity means that the same navigational mathematics applies at every level of description. Whether one is navigating the global state of a physical system or resolving a single coordinate of a sub-problem, the algorithm is identical: backward trajectory completion through a ternary hierarchy using the S-entropy metric. This is the formal content of the observation that ``a sentient cow and a human arrive at $g = 9.81$ m/s$^2$ through incommensurable paths''---both paths are geodesics in isometric S-entropy spaces at different (and intrinsically indistinguishable) hierarchical levels.
\end{remark}

\subsection{The Pendulum Revisited}

\begin{example}[Recursive Structure of the Three-State Pendulum]
\label{ex:pendulum_recursive}
Return to the pendulum with three states $\{p_1, p_2, p_3\}$. Its S-entropy embedding is $\Phi(p_i) = (S_k^{(i)}, S_t^{(i)}, S_e^{(i)})$ for $i = 1,2,3$.

Each state $p_i$ has its own sub-S-entropy structure. Consider $p_1$:
\begin{align}
p_1 &\mapsto (S_k^{(1)}, S_t^{(1)}, S_e^{(1)}) \\
S_e^{(1)} &\mapsto (S_{e,k}^{(1)}, S_{e,t}^{(1)}, S_{e,e}^{(1)})
\end{align}

The sub-coordinates of $S_e^{(1)}$ describe:
\begin{itemize}
\item $S_{e,k}^{(1)}$: How much of the constraint structure at $p_1$ is resolved.
\item $S_{e,t}^{(1)}$: The temporal position of constraint accumulation at $p_1$.
\item $S_{e,e}^{(1)}$: The meta-constraint density (constraints on constraints) at $p_1$.
\end{itemize}

By Theorem~\ref{thm:recursive_metric}, the metric on $(S_{e,k}^{(1)}, S_{e,t}^{(1)}, S_{e,e}^{(1)})$ is isometric to the metric on $(S_k, S_t, S_e)$. Navigating within $p_1$'s constraint structure uses the same algorithm as navigating the global pendulum state.

Furthermore, the compositions of 3 act not only on the global period but on each sub-coordinate: $S_t^{(1)}$ can be decomposed as $(1{+}1{+}1)$, $(1{+}2)$, or $(2{+}1)$ in its own temporal domain, generating the same ternary partition structure recursively.

At every node, the fundamental identity holds: resolving $S_{e,k}^{(1)}$ is simultaneously an observation ($\mathcal{O}$), a computation ($\mathcal{C}$), and a physical process ($\mathcal{P}$). The pendulum contains the entire framework in miniature, repeated at every scale.
\end{example}

\begin{figure*}[!htbp]
    \centering
    \includegraphics[width=\textwidth]{figures/embedding_panel3_self_similarity.png}
    \caption{Recursive Self-Similarity of the S-Entropy Metric. (a) Scatter plot of $200$ pairwise distances computed at two distinct hierarchy scales; the tight clustering around the line $y=x$ (navy dashed) demonstrates that distances are preserved under scale transformations, the core property of the Recursive Metric Theorem. (b) Pass rates for two independent self-similarity tests: ordering preservation across scales ($98.9\%$, $176/178$, navy) and metric structure matching ($100.0\%$, $162/162$, coral); both exceed $90\%$, confirming that the metric is scale-invariant to within numerical precision. (c) Three-dimensional scatter of thirty categorical states rendered simultaneously at three hierarchy scales (scale~1 in navy circles, scale~2 in teal triangles, scale~3 in coral squares); the clusters preserve their spatial arrangement across scales, visualizing the isometry between sub-coordinate and global metric structure. (d) Histogram of distance ratios between the two scales; the distribution is tightly concentrated around unity (coral dashed line) with the measured mean (navy line) matching $1.000$ to within noise, establishing that sub-metrics are isometric to the global metric.}
    \label{fig:embedding_panel3_self_similarity}
\end{figure*}


\section{Backward Navigation as Geodesic Flow}
\label{sec:geodesic}

\subsection{Geodesic Equations on the S-Entropy Manifold}

\begin{theorem}[Backward Navigation is Geodesic]
\label{thm:geodesic_nav}
In the S-entropy manifold $([0,1]^3, g)$ with metric~\eqref{eq:s_metric}, the optimal backward trajectory from final state $\mathbf{S}_f$ to the root is the unique geodesic connecting $\mathbf{S}_f$ to $\mathbf{S}_0 = (1, 0, 0)$.
\end{theorem}

\begin{proof}
The geodesic equations for the diagonal metric $g_{ij} = \delta_{ij} / (S_i(1-S_i))$ are:
\begin{equation}
\label{eq:geodesic}
\ddot{S}_i + \frac{1-2S_i}{2S_i(1-S_i)} \dot{S}_i^2 = 0, \quad i \in \{k, t, e\}
\end{equation}
where dots denote derivatives with respect to the affine parameter.

Each equation is independent (diagonal metric), and each has the form of the geodesic equation on $(0,1)$ with the Fisher metric $ds^2 = dS^2 / (S(1-S))$. The geodesics of this one-dimensional metric are well-known~\cite{amari1985differential}: they are the arcsine reparametrizations
\begin{equation}
\label{eq:geodesic_solution}
S_i(\lambda) = \sin^2\!\left(\frac{\pi}{2} \cdot \frac{\lambda - \lambda_0}{\lambda_f - \lambda_0}\right)
\end{equation}
interpolating between boundary values $S_i(\lambda_0)$ and $S_i(\lambda_f)$ with constant speed in the Fisher metric.

The backward trajectory from $\mathbf{S}_f = (S_k^f, S_t^f, S_e^f)$ (a fully resolved state: $S_k^f$ near 0, $S_t^f$ near 1, $S_e^f$ near 1) to $\mathbf{S}_0 = (1, 0, 0)$ (root: maximum uncertainty, no time elapsed, no constraints) is the product of three such arcsine curves:
\begin{equation}
\gamma(\lambda) = \left(\sin^2\!\left(\frac{\pi\lambda}{2}\right), \sin^2\!\left(\frac{\pi(1-\lambda)}{2}\right), \sin^2\!\left(\frac{\pi(1-\lambda)}{2}\right)\right)
\end{equation}

This geodesic is unique because the Fisher metric on $(0,1)$ is complete and has strictly positive curvature, ensuring unique shortest paths between any two interior points~\cite{docarmo1992riemannian}.

\textbf{Optimality}: The geodesic minimizes the arc length:
\begin{equation}
L[\gamma] = \int_0^1 \sqrt{g_{ij} \dot{S}^i \dot{S}^j} \, d\lambda
\end{equation}
among all paths from $\mathbf{S}_f$ to $\mathbf{S}_0$. Since $L[\gamma]$ equals the S-distance, backward navigation along the geodesic minimizes the total S-distance traversed. $\square$
\end{proof}

\subsection{Complexity of Geodesic Navigation}

\begin{theorem}[Geodesic Navigation Complexity]
\label{thm:geodesic_complexity}
Backward navigation along the geodesic in S-entropy space achieves $\mathcal{O}(\log_3 N)$ complexity for a ternary hierarchy with $N = 3^k$ states.
\end{theorem}

\begin{proof}
The geodesic from $\mathbf{S}_f$ to $\mathbf{S}_0$ crosses $k$ level boundaries in the $S_k$ coordinate (from $S_k \approx 0$ to $S_k = 1$, passing through $S_k = 1 - j/k$ for $j = k-1, k-2, \ldots, 0$).

At each level boundary, the geodesic determines which of three sub-blocks to enter. This ternary decision is resolved by comparing geodesic distances to the three child nodes, requiring $\mathcal{O}(1)$ distance computations per level.

Total decisions: $k = \log_3 N$.
Cost per decision: $\mathcal{O}(1)$ (distance comparison in the continuous metric).
Total complexity: $\mathcal{O}(k) = \mathcal{O}(\log_3 N)$.

This matches the information-theoretic lower bound of $\log_3 N$ ternary decisions needed to identify one leaf among $N$. The geodesic achieves this bound because the navigation-compatible metric ensures that each distance comparison resolves exactly one level of the hierarchy. $\square$
\end{proof}

\subsection{Comparison: Forward vs.\ Backward Complexity}

\begin{center}
\begin{tabular}{lcc}
\toprule
Method & Metric Required? & Complexity \\
\midrule
Forward observation & No & $\mathcal{O}(N)$ \\
Discrete backward search & No & $\Omega(N)$ \\
Binary backward (with metric) & Yes & $\mathcal{O}(\log_2 N)$ \\
S-entropy geodesic (ternary) & Yes & $\mathcal{O}(\log_3 N)$ \\
\midrule
Ratio (S-entropy / forward) & --- & $\log_3 N / N$ \\
\bottomrule
\end{tabular}
\end{center}

The S-entropy geodesic achieves a factor of $N / \log_3 N$ improvement over both forward observation and discrete backward search. For $N = 3^{10} = 59{,}049$, this is a $5{,}905\times$ speedup.

\begin{figure*}[!htbp]
    \centering
    \includegraphics[width=\textwidth]{figures/embedding_panel2_compatibility.png}
    \caption{Navigation Compatibility of the S-Entropy Metric. (a) Pass rates for the three navigation-compatibility conditions across 500 random tests each: refinement monotonicity ($1.000$, navy), sibling separation ($1.000$, teal), and contraction ($0.977$, coral); the dashed line at 0.9 marks the conventional threshold, with all three conditions satisfied. (b) Distribution of child-to-parent diameter ratios across $1{,}500$ tested (parent, child) pairs; the mean ratio is $0.856$ (navy line), and the overwhelming majority lie below unity (coral dashed line), confirming that the metric contracts block diameters at every level of the hierarchy. (c) Three-dimensional surface of block diameter over $(\text{depth},\, \text{sibling index})$; the monotone decrease along the depth axis is the geometric signature of refinement monotonicity, while the symmetric variation along the sibling axis reflects the $\mathbb{Z}_3$ permutation invariance of the ternary hierarchy. (d) Block diameter vs.\ depth on a logarithmic scale; the exponential decay $d \propto 0.856^{\text{depth}}$ is the quantitative content of the contraction property, with the filled region showing the cumulative volume of resolved categorical space at each depth.}
    \label{fig:embedding_panel2_compatibility}
\end{figure*}


\section{Physical Realization}
\label{sec:physical}

\subsection{From Pendulum to General Oscillator}

The pendulum example generalizes to any bounded oscillator.

\begin{theorem}[Universal S-Entropy Embedding]
\label{thm:universal}
Any bounded physical system with $M$ distinguishable modes and $n$ energy levels per mode admits an S-entropy embedding $\Phi: \mathcal{S}_{\text{phys}} \to [0,1]^3$ that is navigation-compatible and satisfies the triple equivalence.
\end{theorem}

\begin{proof}
By the Triple Equivalence Theorem, the system admits three descriptions with identical entropy $S = k_B M \ln n$. The proof of Theorem~\ref{thm:three_coordinates} shows that these three descriptions generate three independent coordinates. The proof of Theorem~\ref{thm:uniqueness} shows that the embedding is unique up to isometry.

The total number of states is $n^M$. In the ternary hierarchy, the depth is $k = \lceil \log_3(n^M) \rceil = \lceil M \log_3 n \rceil$. The S-entropy coordinates are:
\begin{align}
S_k &= 1 - \frac{\text{resolved levels}}{k} \in [0,1] \\
S_t &= \frac{\text{completed states}}{n^M} \in [0,1] \\
S_e &= \frac{\text{accumulated constraints}}{\binom{M}{2}} \in [0,1]
\end{align}

The embedding is navigation-compatible by Proposition~\ref{prop:nav_compat}, and the metric is the Fisher information metric on each coordinate (Theorem~\ref{thm:uniqueness}). $\square$
\end{proof}

\subsection{Processor-Oscillator Duality}

\begin{corollary}[Every Computation Admits S-Entropy Navigation]
\label{cor:computation}
By the processor-oscillator duality---every processor is an oscillator and every oscillator is a processor---every physical computation admits an S-entropy embedding. Backward navigation through the computation's state space is therefore always possible in $\mathcal{O}(\log_3 N)$ time, where $N$ is the number of distinguishable computational states.
\end{corollary}

\begin{proof}
A digital processor with clock frequency $f$ and word length $w$ has $n^w$ distinguishable states (with $n$ values per register). By the processor-oscillator duality, this is identical to an oscillator with $M = w$ modes and $n$ energy levels. Theorem~\ref{thm:universal} provides the S-entropy embedding. $\square$
\end{proof}

\subsection{The Cow and the Physicist}

\begin{example}[Substrate-Independent Convergence]
\label{ex:cow}
Consider two observers measuring gravitational acceleration $g$:
\begin{enumerate}
\item A human physicist with apparatus, using Newtonian mechanics.
\item A sentient cow with apparatus, using an unknown reasoning process.
\end{enumerate}

Both arrive at $g = 9.81$ m/s$^2$. The physicist's trajectory through S-entropy space is a geodesic $\gamma_{\text{human}}(\lambda)$. The cow's trajectory is a different geodesic $\gamma_{\text{cow}}(\lambda)$.

By the universal embedding theorem (Theorem~\ref{thm:universal}), both trajectories lie in the same S-entropy space (since the physical system---Earth's gravity---is the same). By scale ambiguity (Corollary~\ref{cor:no_top}), the cow's internal decomposition of the problem into sub-S-coordinates is isometric to the human's, but at a different (and intrinsically indistinguishable) hierarchical level.

The endpoint is shared: $\mathbf{S}_f = (0, 1, 1)$ (fully resolved, temporally complete, fully constrained). The paths are different: $\gamma_{\text{human}} \neq \gamma_{\text{cow}}$ (different intermediate states, different decompositions, different sub-hierarchies). But both are geodesics in the same metric space, so both achieve $\mathcal{O}(\log_3 N)$ complexity.

The human cannot reconstruct the cow's path from the shared endpoint, because scale ambiguity prevents determining which hierarchical level the cow's intermediate states occupy. The process happened, the answer is correct, the intermediate path is opaque---not because information was lost, but because the S-entropy metric is isometric across levels, making inter-level distinctions unmeasurable.
\end{example}

\begin{theorem}[Path Opacity]
\label{thm:path_opacity}
Let $\gamma_1, \gamma_2$ be two geodesics in S-entropy space sharing the same endpoint $\mathbf{S}_f$ but traversing different hierarchical decompositions. Given only $\mathbf{S}_f$ and the S-entropy metric $g$, it is impossible to determine which geodesic was followed.
\end{theorem}

\begin{proof}
By Corollary~\ref{cor:no_top}, the hierarchical level of any intermediate state is intrinsically unmeasurable. Two geodesics sharing an endpoint but passing through different levels produce the same final S-entropy triple. Since the metric is isometric across levels, no distance measurement at the endpoint can distinguish the two paths.

Formally: let $\gamma_1$ pass through intermediate states at levels $\{n_1, n_2, \ldots\}$ and $\gamma_2$ through levels $\{m_1, m_2, \ldots\}$. For any metric invariant $I$ computed from $\mathbf{S}_f$ and $g$:
\begin{equation}
I(\gamma_1, \mathbf{S}_f, g) = I(\gamma_2, \mathbf{S}_f, g)
\end{equation}
because $I$ depends only on the metric structure at $\mathbf{S}_f$, which is level-independent.

Therefore the path is opaque: the endpoint contains no record of which geodesic was followed. $\square$
\end{proof}


\section{The $3 \times 3$ Structure}
\label{sec:3x3}

\subsection{Three Layers of Triple Equivalence}

The framework exhibits a $3 \times 3$ structure: three layers of description, each containing three equivalent elements.

\begin{definition}[The Triple-of-Triples]
\label{def:triple_of_triples}
The complete structural content of the S-entropy framework is the $3 \times 3$ matrix:
\begin{equation}
\label{eq:3x3}
\mathcal{F} = \begin{pmatrix}
\text{Oscillation} & \text{Partition} & \text{Category} \\
\text{Observation} & \text{Computation} & \text{Processing} \\
S_t & S_k & S_e
\end{pmatrix}
\end{equation}
where:
\begin{itemize}
\item Row 1 (Description): what the system \emph{is}---three equivalent descriptions.
\item Row 2 (Operation): what the system \emph{does}---three equivalent operations.
\item Row 3 (Coordinate): where the system \emph{is}---three equivalent measurements.
\end{itemize}
\end{definition}

\begin{theorem}[Row and Column Equivalences]
\label{thm:row_column}
In the $3 \times 3$ structure $\mathcal{F}$:
\begin{enumerate}
\item Each \textbf{row} is a triple equivalence (already established).
\item Each \textbf{column} is a triple equivalence.
\end{enumerate}
\end{theorem}

\begin{proof}
Row equivalences are the established results: Description ($S_O = S_C = S_P$), Operation ($\mathcal{O} \equiv \mathcal{C} \equiv \mathcal{P}$), and Coordinate ($S_k, S_t, S_e$ are projections of a single state).

For column equivalences, we must show that each column's three elements are structurally identical:

\textbf{Column 1 (Temporal)}: Oscillation--Observation--$S_t$.
\begin{itemize}
\item Oscillation is temporal: period $T = 2\pi/\omega$, frequency $f$, phase $\phi(t)$.
\item Observation is a temporal act: it occurs at a specific time and has a duration.
\item $S_t$ measures temporal progress through the completion sequence.
\end{itemize}
All three are parameterized by a single variable (time), evolving monotonically, and quantifying the same temporal structure from different perspectives. By the fundamental identity, observing ($\mathcal{O}$) the oscillation phase ($\Osc$) at time $S_t$ are three descriptions of a single temporal event.

\textbf{Column 2 (Structural)}: Partition--Computation--$S_k$.
\begin{itemize}
\item Partition is structural: block decomposition, refinement depth, equivalence classes.
\item Computation is structural: information transformation, input-output mapping, algorithm.
\item $S_k$ measures structural resolution---the fraction of the partition hierarchy resolved.
\end{itemize}
All three describe the informational structure of the system. Partitioning ($\Part$), computing ($\mathcal{C}$), and measuring $S_k$ are three descriptions of the same structural decomposition.

\textbf{Column 3 (Dynamical)}: Category--Processing--$S_e$.
\begin{itemize}
\item Category is dynamical: morphisms between states, endomorphisms, state transitions.
\item Processing is dynamical: physical evolution, energy exchange, state change.
\item $S_e$ measures dynamical constraint density---how constrained the transitions are.
\end{itemize}
All three describe the dynamical aspect: which transitions are possible, at what cost, under what constraints. Categorical morphism ($\Cat$), physical processing ($\mathcal{P}$), and constraint measurement ($S_e$) are three descriptions of the same dynamical structure.

Each column equivalence follows from the fundamental identity applied to the corresponding aspect (temporal, structural, dynamical) of the system. $\square$
\end{proof}

\begin{corollary}[Nine Names, Three Things, One Object]
\label{cor:nine_three_one}
The $3 \times 3$ structure $\mathcal{F}$ has nine entries but only three independent elements (one per column, or equivalently one per row). By the triple equivalence in each row and each column, all nine entries are equivalent up to perspective. The entire framework reduces to a single object viewed from nine equally valid perspectives.
\end{corollary}

\begin{proof}
Each row collapses three entries to one (by row equivalence). Each column collapses three entries to one (by column equivalence). Since rows and columns intersect, the three row-equivalence-classes are themselves equivalent (each contains one element from each column, and column equivalence identifies these). Therefore all nine entries are equivalent. $\square$
\end{proof}

\subsection{Self-Similarity of the $3 \times 3$ Structure}

\begin{theorem}[Fractal $3 \times 3$]
\label{thm:fractal_3x3}
Each entry of $\mathcal{F}$ is itself a $3 \times 3$ structure isomorphic to $\mathcal{F}$.
\end{theorem}

\begin{proof}
By Theorem~\ref{thm:recursive_metric}, each S-entropy coordinate decomposes into three sub-coordinates with isometric metric. The sub-coordinates inherit the triple equivalence structure (the sub-hierarchy has oscillator, partition, and categorical descriptions). Therefore the sub-structure has its own $3 \times 3$ matrix $\mathcal{F}_{\text{sub}}$, and by the uniqueness theorem (Theorem~\ref{thm:uniqueness}), $\mathcal{F}_{\text{sub}}$ is isomorphic to $\mathcal{F}$.

This applies to every entry, generating an infinite descent:
\begin{equation}
\mathcal{F} \to \mathcal{F}_{i,j} \to \mathcal{F}_{i,j,k,l} \to \cdots
\end{equation}
with $\mathcal{F}_{i,j} \cong \mathcal{F}$ at every level. $\square$
\end{proof}


\section{Virtual Sub-States and Non-Physical Intermediate Decompositions}
\label{sec:virtual}

\subsection{The Miracle Theorem}

Backward navigation through S-entropy space has a property with no analogue in forward computation: the intermediate sub-task decompositions need not be individually physical.

\begin{theorem}[Virtual Sub-State Theorem]
\label{thm:virtual}
Let $\mathbf{S} = (S_k, S_t, S_e) \in [0,1]^3$ be a valid global S-entropy state. Let $S_i$ for $i \in \{k, t, e\}$ decompose into sub-coordinates $\mathbf{S}_i = (S_{i,k}, S_{i,t}, S_{i,e})$ via Theorem~\ref{thm:recursive_metric}. Then:
\begin{enumerate}
\item The sub-coordinates $S_{i,j}$ are not individually constrained to $[0,1]$.
\item There exist valid global states $\mathbf{S} \in [0,1]^3$ whose sub-decompositions contain $S_{i,j} < 0$ (negative values) or $S_{i,j} > 1$ (super-unit values).
\item The global state $\mathbf{S}$ remains physically valid despite non-physical sub-coordinates, provided the composition map $f: \mathbf{S}_i \mapsto S_i$ yields $S_i \in [0,1]$.
\end{enumerate}
\end{theorem}

\begin{proof}
\textbf{Part 1: Sub-coordinates are unconstrained.}

The global S-entropy coordinates are constrained to $[0,1]^3$ by their physical definitions (Definitions~\ref{def:s_k}--\ref{def:s_e}): $S_k$ is a ratio of resolved levels, $S_t$ is a ratio of completed states, $S_e$ is a ratio of accumulated constraints. Each ratio is bounded by construction.

The sub-coordinates $S_{i,j}$ arise from the recursive decomposition (Theorem~\ref{thm:recursive_metric}). The decomposition is defined by the requirement that the sub-metric be isometric to the global metric---not by the requirement that sub-coordinates satisfy independent physical bounds.

Formally: the isometry $\iota: ([0,1]^3, g) \to ([0,1]^3, g|_{\mathbf{S}_i})$ maps the global coordinate system to the sub-coordinate system. But $\iota$ is a metric isometry, not a set inclusion. The sub-coordinates $S_{i,j}$ are defined as the image of the ternary hierarchy under the sub-embedding $\Phi_i$, and $\Phi_i$ maps into the tangent space of $[0,1]^3$ at the point $S_i$---a space that is locally $\mathbb{R}^3$, not $[0,1]^3$.

Therefore sub-coordinates range over $\mathbb{R}$, not $[0,1]$.

\textbf{Part 2: Non-physical sub-coordinates exist.}

Consider a global state $\mathbf{S} = (0.5, 0.5, 0.5)$ (midpoint of the cube). The knowledge coordinate $S_k = 0.5$ decomposes into sub-coordinates $(S_{k,k}, S_{k,t}, S_{k,e})$ describing the resolution process.

In backward navigation, the resolution of $S_k$ proceeds from the known final state backward. At intermediate steps, the navigator may ``overshoot'' in the temporal sub-coordinate: $S_{k,t} > 1$ (the sub-task accesses temporal information beyond the endpoint) or ``borrow'' from the knowledge sub-coordinate: $S_{k,k} < 0$ (the sub-task uses knowledge that has not yet been acquired at that resolution level).

These overshoots are permissible because the composition $f(S_{k,k}, S_{k,t}, S_{k,e}) = S_k = 0.5$ remains in $[0,1]$. The overshoots cancel in the composition.

Concretely: let $f(\mathbf{S}_k) = \frac{1}{3}(S_{k,k} + S_{k,t} + S_{k,e})$ (equal-weight composition). Then the decomposition $(S_{k,k}, S_{k,t}, S_{k,e}) = (-0.2, 0.8, 0.9)$ yields $f = 0.5 \in [0,1]$ despite $S_{k,k} = -0.2 < 0$.

\textbf{Part 3: Physical validity of the global state.}

The global state $\mathbf{S} \in [0,1]^3$ is the physically observable quantity. By the path opacity theorem (Theorem~\ref{thm:path_opacity}), the sub-coordinate decomposition is unobservable from the global state. Therefore the non-physicality of individual sub-coordinates has no observable consequence: no measurement at the global level can detect whether the sub-decomposition passed through non-physical values.

The global state is valid if and only if $\mathbf{S} \in [0,1]^3$. The sub-coordinates are virtual---they participate in the navigation algorithm but are not independently measurable. $\square$
\end{proof}

\begin{definition}[Virtual Sub-State]
\label{def:virtual}
A \emph{virtual sub-state} is a sub-coordinate triple $\mathbf{S}_i = (S_{i,k}, S_{i,t}, S_{i,e})$ with at least one component outside $[0,1]$, occurring as an intermediate decomposition of a valid global state $\mathbf{S} \in [0,1]^3$.
\end{definition}

\subsection{Analogy with Virtual Particles}

\begin{remark}[Quantum Field Theory Parallel]
\label{rem:qft}
The virtual sub-state phenomenon is structurally identical to virtual particles in quantum field theory~\cite{hatfield1992quantum}. In QFT:
\begin{itemize}
\item A real process (e.g., electron-electron scattering) has physical initial and final states satisfying energy-momentum conservation.
\item Intermediate states (virtual photon exchange) violate energy conservation: $\Delta E \cdot \Delta t < \hbar/2$.
\item The violation is unobservable: only the initial and final states are measured.
\item The ``impossible'' intermediate states are essential to the calculation---removing them gives wrong answers.
\end{itemize}

In S-entropy:
\begin{itemize}
\item A real navigation (backward trajectory completion) has physical initial and final S-values in $[0,1]^3$.
\item Intermediate sub-decompositions may violate $S_{i,j} \in [0,1]$: negative time, super-unit knowledge, negative constraint density.
\item The violation is unobservable: only the global S-values are measured (path opacity).
\item The ``impossible'' sub-states are essential to the navigation---restricting to physical sub-states reduces navigational flexibility and can increase complexity.
\end{itemize}

The parallel extends to the uncertainty principle. In QFT, $\Delta E \cdot \Delta t < \hbar/2$ bounds the product of energy violation and duration. In S-entropy, the analogous bound is:
\begin{equation}
\label{eq:s_uncertainty}
|S_{i,j} - S_{i,j}^{\text{phys}}| \cdot \Delta_{\text{level}} < \epsilon
\end{equation}
where $S_{i,j}^{\text{phys}}$ is the nearest physical value, $\Delta_{\text{level}}$ is the hierarchical depth of the sub-coordinate, and $\epsilon$ bounds the total deviation. Deeper sub-coordinates (larger $\Delta_{\text{level}}$) are permitted smaller deviations from physicality, while shallow sub-coordinates may deviate more---but the product remains bounded.
\end{remark}

\subsection{Why Forward Navigation Cannot Use Virtual Sub-States}

\begin{theorem}[Asymmetry of Virtual States]
\label{thm:forward_no_virtual}
Forward navigation (from initial state to final state by step-wise construction) cannot use virtual sub-states. Only backward navigation (from known final state) permits non-physical intermediate decompositions.
\end{theorem}

\begin{proof}
In forward navigation, each step constructs a new state from the current state. The new state must be physically realizable because it becomes the input to the next step. If an intermediate state has $S_{i,j} < 0$, the next step receives a non-physical input and cannot proceed---there is no physical process that starts from a negative-time configuration.

Formally: forward navigation composes functions $f_1, f_2, \ldots, f_k$ sequentially:
\begin{equation}
\mathbf{S}_0 \xrightarrow{f_1} \mathbf{S}_1 \xrightarrow{f_2} \cdots \xrightarrow{f_k} \mathbf{S}_f
\end{equation}
Each $f_j$ must have domain containing physical states ($[0,1]^3$) and range within physical states. Therefore every intermediate $\mathbf{S}_j \in [0,1]^3$.

In backward navigation, the trajectory is not constructed step-by-step but recovered from the known endpoint. The sub-decompositions are not inputs to physical processes; they are components of the backward analysis. The navigation algorithm evaluates the geodesic equation (Equation~\ref{eq:geodesic}) at sub-coordinate resolution, and the geodesic passes through the full tangent space $\mathbb{R}^3$, not just $[0,1]^3$.

The asymmetry is fundamental: forward navigation is constructive (each step must be physically executable), while backward navigation is analytical (the endpoint is given, and intermediate decompositions are analytical tools). Virtual sub-states are analytical artefacts of the backward decomposition, not physical states that must be realized. $\square$
\end{proof}

\begin{figure*}[!htbp]
    \centering
    \includegraphics[width=\textwidth]{figures/embedding_panel4_miracles.png}
    \caption{Miracle Resolution Along the Backward Trajectory. (a) Miracle count trajectory for a depth-10 hierarchy; the count decreases monotonically from 10 (all levels unresolved, maximum miracle budget) to 1 (penultimate state, single completion morphism remaining, coral dashed line). The filled region quantifies the cumulative miracle-resolution work along the backward geodesic. (b) Miracle count trajectories for all eight tested hierarchy depths ($d=3$ through $d=10$, coloured by depth via viridis); every trajectory terminates at miracle count~1, confirming that the minimum miracle principle identifies the penultimate state independently of hierarchy size. (c) Three-dimensional surface of miracle count over $(\text{backward step},\, \text{hierarchy depth})$; the monotone descent along the step axis and linear growth along the depth axis together visualize the claim that miracle count equals the number of remaining unresolved ternary decisions. (d) Final miracle count at trajectory termination vs.\ hierarchy depth; all eight depths yield final count $=1$ (coral dashed reference), confirming that the minimum miracle principle identifies the penultimate state uniformly across scales.}
    \label{fig:embedding_panel4_miracles}
\end{figure*}

\subsection{Consequence: Backward Navigation is Strictly More Powerful}

\begin{corollary}[Strict Power Separation]
\label{cor:strict_power}
The class of problems solvable by backward navigation in S-entropy space strictly contains the class solvable by forward navigation, because backward navigation can traverse regions of sub-S-space inaccessible to forward navigation.
\end{corollary}

\begin{proof}
Forward navigation is confined to $[0,1]^3$ at every intermediate step. Backward navigation, through virtual sub-states, accesses the full tangent bundle $T([0,1]^3) \cong [0,1]^3 \times \mathbb{R}^3$.

The geodesic connecting two points in $[0,1]^3$ may pass through regions of sub-coordinate space outside $[0,1]$. If the only geodesic connecting $\mathbf{S}_0$ to $\mathbf{S}_f$ passes through virtual sub-states, then forward navigation (constrained to $[0,1]^3$) cannot follow this geodesic and must find an alternative path---which may be longer (higher complexity) or may not exist at all.

Therefore backward navigation solves all problems solvable by forward navigation (by restricting to physical sub-states) and additionally solves problems requiring virtual sub-state traversal. The inclusion is strict whenever the optimal geodesic passes through virtual territory. $\square$
\end{proof}

\subsection{Virtual Sub-States Are Necessary for the Complexity Advantage}

The preceding results show that virtual sub-states are \emph{permitted}. We now prove the stronger claim: they are \emph{necessary}. Without them, backward navigation loses its complexity advantage entirely.

\begin{theorem}[Collapse Without Virtual States]
\label{thm:collapse}
Backward navigation in S-entropy space, restricted to physically realizable sub-coordinate decompositions ($S_{i,j} \in [0,1]$ at every level of the recursion), has worst-case complexity $\Omega(N)$---identical to forward simulation.
\end{theorem}

\begin{proof}
\textbf{Step 1: Physically-restricted backward navigation is time-reversed forward simulation.}

Forward simulation from initial state $\mathbf{S}_0$ to final state $\mathbf{S}_f$ applies a sequence of physically realizable transitions:
\begin{equation}
\mathbf{S}_0 \xrightarrow{f_1} \mathbf{S}_1 \xrightarrow{f_2} \cdots \xrightarrow{f_k} \mathbf{S}_f
\end{equation}
where each $f_j$ is determined by the equations of motion (physical laws) and each $\mathbf{S}_j \in [0,1]^3$.

Backward navigation from $\mathbf{S}_f$, restricted to physical sub-coordinates, must find the sequence in reverse:
\begin{equation}
\mathbf{S}_f \xrightarrow{f_k^{-1}} \mathbf{S}_{k-1} \xrightarrow{f_{k-1}^{-1}} \cdots \xrightarrow{f_1^{-1}} \mathbf{S}_0
\end{equation}
where each $f_j^{-1}$ must produce a physical state $\mathbf{S}_{j-1} \in [0,1]^3$.

To compute $f_j^{-1}(\mathbf{S}_j)$, the algorithm must know which physical law $f_j$ was applied at step $j$. This requires either:
\begin{enumerate}
\item[(a)] Knowing the equations of motion a priori (equivalent to forward simulation).
\item[(b)] Discovering the equations of motion from the state $\mathbf{S}_j$ (requires $\Omega(N)$ queries by Theorem~\ref{thm:discrete_lower}, since the physical law is encoded in the hierarchy structure).
\end{enumerate}

In either case, each backward step costs at least as much as one forward step. Over $k$ steps, the total cost is $\Omega(k)$---but determining \emph{which} forward step to invert at each level requires identifying the correct branch of the ternary hierarchy, which without a metric shortcut costs $\Omega(N/3^j)$ at level $j$. Summing:
\begin{equation}
\sum_{j=0}^{k-1} \frac{N}{3^j} = N \cdot \frac{1 - 3^{-k}}{1 - 3^{-1}} = \frac{3}{2}(N - 1) = \Omega(N)
\end{equation}

\textbf{Step 2: The geodesic shortcut requires virtual territory.}

Geodesic navigation achieves $\mathcal{O}(\log_3 N)$ by cutting through S-entropy space along the shortest path. At each level, the geodesic resolves one ternary decision by distance comparison rather than physical-law inversion.

The distance comparison evaluates $d(\mathbf{S}_f, c_i)$ for the three child centroids $c_0, c_1, c_2$ (Theorem~\ref{thm:metric_nav}). This evaluation involves computing the geodesic distance, which requires evaluating the metric along the path---including at sub-coordinate resolution where the path passes through the tangent space $\mathbb{R}^3$.

If the sub-coordinates are restricted to $[0,1]$, the geodesic distance computation is constrained to the interior of the cube. But the shortest geodesic between two points in $[0,1]^3$ under the Fisher metric may exit the cube at the sub-coordinate level (the Fisher metric has positive curvature, and geodesics on positively curved spaces can ``bulge'' outside coordinate bounds at finer resolution).

Restricting to the interior forces the algorithm to follow a longer path---one that avoids virtual sub-coordinate territory. This longer path corresponds exactly to the step-by-step physical trajectory, recovering the $\Omega(N)$ cost.

\textbf{Step 3: Conclusion.}

The $\mathcal{O}(\log_3 N)$ complexity of backward navigation is achieved by the geodesic shortcut. The geodesic shortcut traverses virtual sub-state territory. Removing virtual sub-states removes the shortcut, collapsing backward navigation to time-reversed forward simulation at $\Omega(N)$ cost.

Therefore virtual sub-states are not merely permitted but \emph{necessary} for the complexity advantage. $\square$
\end{proof}

\begin{corollary}[Virtual States Are the Source of Speedup]
\label{cor:speedup_source}
The speedup factor $N / \log_3 N$ of backward navigation over forward simulation is entirely attributable to virtual sub-states. Specifically, the fraction of the optimal geodesic that passes through virtual sub-coordinate territory is:
\begin{equation}
\frac{L_{\text{virtual}}}{L_{\text{total}}} = 1 - \frac{\log_3 N}{N} \to 1 \text{ as } N \to \infty
\end{equation}
In the large-$N$ limit, essentially the entire geodesic is virtual.
\end{corollary}

\begin{remark}[The Pendulum Without Laws]
\label{rem:pendulum_no_laws}
Consider recovering the trajectory of a pendulum from its final state $p_3$.

\emph{Forward simulation}: requires knowing the equation of motion $\ddot{\theta} + (g/L)\sin\theta = 0$, initial conditions $(\theta_0, \dot{\theta}_0)$, and numerical integration. Each step applies the law and produces a physical intermediate state.

\emph{Physically-restricted backward navigation}: requires inverting the equation of motion at each step. To find $p_2$ from $p_3$, one must know what force produced the transition $p_2 \to p_3$ and invert it. This is forward simulation run backward---same equations, same cost.

\emph{Backward navigation with virtual sub-states}: requires only the S-entropy coordinates of $p_3$ and the geodesic equation~\eqref{eq:geodesic}. The geodesic determines $p_2$ by distance minimization in S-entropy space, without reference to the equation of motion. The sub-coordinate decomposition of the geodesic at the $p_2 \to p_3$ transition may involve $S_{k,t} < 0$ (the ``temporal sub-task'' accesses a moment before the pendulum existed) or $S_{k,k} > 1$ (the ``knowledge sub-task'' knows more than is physically available at $p_2$). These virtual sub-states are how the geodesic bypasses the need for the equation of motion.

No physical law is applied. No equation is inverted. The step before $p_1$ is found by navigating S-entropy space through territory that no physical pendulum ever occupies. This is why some sub-states \emph{must} be impossible: the backward trajectory does not follow the physical trajectory, and the regions it does traverse have no physical counterpart.
\end{remark}

\begin{remark}[The Precise Meaning of ``Miracle'']
\label{rem:miracle}
A \emph{miracle} in the S-entropy framework is a backward trajectory whose sub-task decomposition contains virtual sub-states. The trajectory happened (the global S-value is valid). The intermediate sub-tasks are individually non-physical (negative time, borrowed knowledge, impossible constraints). No reconstruction of the intermediate path is possible from the endpoint (path opacity). The event is real, the outcome is correct, the explanation is structurally inaccessible.

This is not metaphor. It is a theorem: Theorem~\ref{thm:virtual} proves virtual sub-states exist, Theorem~\ref{thm:path_opacity} proves they are unobservable, Theorem~\ref{thm:forward_no_virtual} proves they are exclusive to backward navigation, and Theorem~\ref{thm:collapse} proves they are necessary for the complexity advantage. The conjunction---real outcome, impossible intermediates, opaque path, necessary for efficiency---is the mathematical content of the word ``miracle.''
\end{remark}

\subsection{Virtual Sub-States as Maxwell Demons}

The virtual sub-state phenomenon has a precise physical identity: each virtual sub-state is a Maxwell demon.

\begin{definition}[Zero-Transfer Transition]
\label{def:zero_transfer}
A \emph{zero-transfer transition} is a state transition $\mathbf{S}_j \to \mathbf{S}_{j+1}$ in which no information about the transition mechanism is transferred to or retained by the final state. That is, from $\mathbf{S}_{j+1}$ alone, it is impossible to determine which mechanism produced the transition.
\end{definition}

\begin{theorem}[Virtual Sub-States Are Maxwell Demons]
\label{thm:virtual_demons}
Each virtual sub-state in a backward geodesic trajectory is a Maxwell demon in the sense of Mizraji~\cite{mizraji2021biological}: an information catalyst that enables a state transition with zero information transfer about the catalytic mechanism.
\end{theorem}

\begin{proof}
We verify the four defining properties of a Maxwell demon against the properties of virtual sub-states.

\textbf{Property 1: Probability transformation.}

A Maxwell demon transforms a near-zero transition probability into a near-unity probability~\cite{mizraji2021biological}. In forward navigation, the probability of selecting the correct ternary branch at depth $j$ without a metric is $1/3^{k-j}$ (uniform random guess over remaining leaves). With the virtual sub-state mediating the geodesic distance comparison, the probability becomes 1 (the correct branch is identified with certainty by distance minimization).

The probability ratio:
\begin{equation}
\frac{p_{\text{with virtual state}}}{p_{\text{without}}} = 3^{k-j}
\end{equation}
which ranges from 3 (at the penultimate level) to $N = 3^k$ (at the root). This is the characteristic $10^6$--$10^{11}$ probability enhancement of biological Maxwell demons.

\textbf{Property 2: Not consumed.}

An enzyme is not consumed by the reaction it catalyses; it appears unchanged in the net equation. A virtual sub-state is not physically realized---it exists only as an intermediate in the geodesic computation, not as a persistent state. It does not appear in the global S-value (by the composition property: the virtual sub-coordinates compose back to a physical parent). Like the enzyme, the virtual sub-state enables the transition without being consumed by it.

\textbf{Property 3: Zero information transfer about mechanism.}

By the path opacity theorem (Theorem~\ref{thm:path_opacity}), the final state $\mathbf{S}_f$ contains no record of which virtual sub-states mediated the backward trajectory. The \emph{outcome} of each ternary decision is recorded (as the categorical address), but the \emph{mechanism} (which virtual sub-coordinate values were used to make the distance comparison) is lost. This is zero-transfer: the information that selected the correct branch is not retained in the final state.

This is precisely the structure of Maxwell's demon: the demon measures which side of the partition a molecule is on (acquiring information), opens the door (enabling the transition), and then erases its memory (by Landauer's principle~\cite{landauer1961irreversibility,bennett1982thermodynamics}). The demon's measurement information does not appear in the final state of the gas---only the outcome (molecule sorted) is retained.

\textbf{Property 4: Coupled input-output filtering.}

From Mizraji's formulation~\cite{mizraji2021biological}, a Maxwell demon operates through coupled filters $\Im_{\text{input}} \circ \Im_{\text{output}}$. In the virtual sub-state:
\begin{itemize}
\item $\Im_{\text{input}}$: the sub-coordinate decomposition of the current state $\mathbf{S}_j$. This decomposes the physical state into (possibly non-physical) sub-coordinates, filtering the vast space of possible decompositions to the specific one used by the geodesic.
\item $\Im_{\text{output}}$: the distance comparison among three child centroids. This filters the three possible next states to the unique correct one.
\item The coupling: $\Im_{\text{input}}$ determines which sub-coordinate values are available, and these values determine which child minimizes the geodesic distance (via $\Im_{\text{output}}$).
\end{itemize}

All four properties are satisfied. Virtual sub-states are Maxwell demons. $\square$
\end{proof}

\begin{corollary}[The Enzyme Analogy Is Exact]
\label{cor:enzyme}
The analogy between virtual sub-states and enzymes is not metaphorical but structural:
\begin{center}
\begin{tabular}{lll}
\toprule
Property & Enzyme & Virtual sub-state \\
\midrule
Enables transition & Chemical reaction & Geodesic step \\
Not consumed & Unchanged after catalysis & Not physically realized \\
Not in net equation & Absent from products & Absent from global $\mathbf{S}$ \\
Lowers barrier & Activation energy $\downarrow$ & Complexity $N \to \log_3 N$ \\
Probability transform & $p_0 \to p_{\text{BMD}} \gg p_0$ & $1/3^k \to 1$ \\
Zero mechanism transfer & Substrate doesn't ``know'' enzyme & $\mathbf{S}_f$ doesn't ``know'' sub-states \\
\bottomrule
\end{tabular}
\end{center}
Both are instances of the same mathematical object: a categorical aperture operating through geometric selection with zero information cost.
\end{corollary}

\begin{figure*}[!htbp]
    \centering
    \includegraphics[width=\textwidth]{figures/embedding_panel5_complexity.png}
    \caption{Forward vs.\ Backward Complexity Scaling. (a) Step counts vs.\ $N$ (log-log) for forward enumeration ($\mathcal{O}(N)$, coral) and backward geodesic navigation ($\mathcal{O}(\log_3 N)$, navy) across twelve hierarchy depths ($N$ ranging from 9 to $1.59\times 10^6$); the asymptotic divergence between the two curves quantifies the complexity advantage of backward navigation. (b) Backward step count vs.\ $\log_3 N$ with linear fit slope $= 1.000$ and $R^2 = 1.000000$; every hierarchy level requires exactly one backward step, saturating the information-theoretic lower bound. (c) Three-dimensional surface of $\log_{10}$(steps) over $(\log_3 N,\, \text{method})$; the divergence between the forward surface (upper) and backward surface (lower) grows with $\log_3 N$, reflecting the $N/\log_3 N$ speedup ratio. (d) Measured speedup of backward over forward navigation vs.\ $N$ (log-log); the speedup reaches $122{,}640\times$ at $N = 1.59\times 10^6$, tracking the theoretical $N/\log_3 N$ reference (gold dashed line) across all tested scales and confirming that the $\mathcal{O}(\log_3 N)$ complexity is achieved in practice, not only in theory.}
    \label{fig:embedding_panel5_complexity}
\end{figure*}

\subsection{The Category Error Correction: From Demon to Aperture}

The identification of virtual sub-states as Maxwell demons requires a critical refinement.

\begin{theorem}[The Corrected Demon]
\label{thm:corrected_demon}
Maxwell's traditional demon commits a category error: it manipulates velocity (kinetic property) expecting to affect entropy (configurational property), but $\partial S_{\text{config}} / \partial v = 0$ because entropy counts spatial arrangements, not velocities~\cite{maxwell1871theory,landauer1961irreversibility,bennett1982thermodynamics}. Virtual sub-states avoid this error entirely. They operate in the correct ontological category: they manipulate categorical addresses, which are configurational properties---exactly what entropy depends on.
\end{theorem}

\begin{proof}
\textbf{The traditional demon's error}: Entropy is $S = k_B \ln \Omega(N_A, N_B)$ where $\Omega$ is the arrangement count---a purely combinatorial quantity depending on spatial distribution $(N_A, N_B)$, not on velocities. The demon sorts by velocity, a kinetic property orthogonal to the arrangement count: $\partial \Omega / \partial v = 0$. The demon operates in the wrong category.

\textbf{The virtual sub-state's correction}: Virtual sub-states operate on S-entropy coordinates $(S_k, S_t, S_e)$, which are derived from the partition description (Definition~\ref{def:s_k}), the oscillator description (Definition~\ref{def:s_t}), and the categorical description (Definition~\ref{def:s_e}). All three descriptions contribute to the arrangement count $\Omega$ through the Triple Equivalence: $S = k_B M \ln n$. The S-entropy coordinates are configurational---they parameterize positions in the partition hierarchy, not velocities.

Therefore the virtual sub-state manipulates a quantity that entropy actually counts: categorical position. The partial derivative is non-zero:
\begin{equation}
\frac{\partial S}{\partial \alpha} \neq 0
\end{equation}
where $\alpha$ is the categorical address. This is why the virtual sub-state succeeds where the demon fails---it operates on the correct variable.

\textbf{Zero cost as a consequence}: The zero information cost of virtual sub-states is not a clever trick that avoids Landauer erasure. It is a structural consequence of operating in the correct category. The commutation relation $[\hat{O}_{\text{cat}}, \hat{O}_{\text{phys}}] = 0$ means categorical and physical operators can be simultaneously evaluated without mutual disturbance. No measurement of the physical state is required because the categorical aperture selects by geometric compatibility, not by information acquisition. No measurement means no erasure, hence zero Landauer cost. $\square$
\end{proof}

\begin{definition}[Categorical Aperture]
\label{def:aperture}
A \emph{categorical aperture} is a topological structure in S-entropy space that permits passage of states based on configurational compatibility rather than kinetic properties. An aperture at depth $j$ in the ternary hierarchy admits states whose categorical address prefix matches the aperture's address up to depth $j$, and rejects all others. The selection is geometric (address matching), not informational (measurement and decision).
\end{definition}

\begin{theorem}[Five Names, One Object]
\label{thm:five_names}
The following are five descriptions of a single mathematical object, listed from most precise to most approximate:
\begin{enumerate}
\item \textbf{Categorical aperture}: geometric selection by configurational compatibility in S-entropy space. Zero information cost. No category error.
\item \textbf{Virtual sub-state}: a non-physical intermediate in a backward geodesic trajectory. Equivalent to (1) by Theorem~\ref{thm:corrected_demon}.
\item \textbf{Miracle}: an event whose outcome is valid but whose intermediate mechanism is non-physical and unrecoverable. Equivalent to (2) by path opacity.
\item \textbf{Enzyme} (categorical interpretation): a biological system that creates categorical apertures for substrate-product transitions. Equivalent to (1) when catalysis is understood as geometric selection rather than temporal acceleration.
\item \textbf{Maxwell demon} (corrected): an information catalyst operating in the configurational category rather than the kinetic category. Historical approximation---overstates the mechanism by implying measurement and decision where geometric selection suffices.
\end{enumerate}
All five share four properties: (a) enables a transition, (b) is not consumed or retained, (c) transfers zero information about itself to the outcome, and (d) is necessary for the transition to occur at the observed efficiency.
\end{theorem}

\begin{proof}
Properties (a)--(d) for categorical apertures: by Definition~\ref{def:aperture}, an aperture enables transition (a), is a geometric structure not consumed by passage (b), retains no information about which state passed through (c), and is necessary because without the aperture all three children are equidistant in the discrete metric (d).

Properties (a)--(d) for virtual sub-states: Theorems~\ref{thm:virtual},~\ref{thm:forward_no_virtual},~\ref{thm:path_opacity},~\ref{thm:collapse}.

Properties (a)--(d) for miracles: Remark~\ref{rem:miracle}.

Properties (a)--(d) for enzymes: Corollary~\ref{cor:enzyme}, with the refinement that ``lowers barrier'' is more precisely ``creates categorical aperture.''

Properties (a)--(d) for Maxwell demons (corrected): Theorem~\ref{thm:virtual_demons}, with the refinement from Theorem~\ref{thm:corrected_demon} that the demon must operate on categorical addresses, not velocities.

All five satisfy the same four-property characterization. By extensionality, they are the same object, listed at decreasing precision: categorical aperture is the exact description, Maxwell demon is the historical approximation. $\square$
\end{proof}

\begin{remark}[Why the Zero-Cost Property Holds]
The zero thermodynamic cost of categorical operations---the commutation relation $[\hat{O}_{\text{cat}}, \hat{O}_{\text{phys}}] = 0$ implying $W_{\text{sort}} = 0$---is a direct consequence of operating in the correct ontological category. Entropy counts spatial arrangements: $S = k_B \ln \Omega(N_A, N_B)$, where $\Omega$ depends on configurational assignments, not on velocities. Categorical operations manipulate configurational assignments (addresses in the partition hierarchy). Physical operations manipulate kinetic properties (velocities, momenta). These are orthogonal categories---the commutation relation $[\hat{O}_{\text{cat}}, \hat{O}_{\text{phys}}] = 0$ is the operator-algebraic expression of this orthogonality. The traditional Maxwell demon pays Landauer cost because it must \emph{measure} a kinetic property (velocity), \emph{decide} based on the measurement, and \emph{erase} the measurement record. The categorical aperture pays zero cost because it performs no measurement---it selects by geometric compatibility, which is a configurational operation on the variable that entropy actually counts.
\end{remark}

\begin{remark}[Closing the Loop]
This identification unifies four papers into one framework:
\begin{itemize}
\item The \emph{trajectory completion} paper establishes backward navigation at $\mathcal{O}(\log_3 N)$ complexity.
\item The \emph{continuous embedding} paper (this paper) identifies the mechanism: geodesic flow through virtual sub-states in S-entropy space.
\item The \emph{St-Stellas} paper identifies the virtual sub-states as biological information catalysts operating through categorical equivalence classes.
\item The \emph{Maxwell demon resolution} paper proves that entropy counts configurations not velocities, establishing why categorical operations (which manipulate configurations) cost zero while kinetic operations (which manipulate velocities) pay Landauer cost.
\end{itemize}
The four papers describe the same phenomenon from four perspectives---description, dynamics, mechanism, and thermodynamic foundation---which is itself an instance of the framework's self-similar structure. All four yield the same conclusion: backward trajectory completion through bounded phase space, mediated by categorical apertures that select configurations at zero information cost, achieving superlinear speedup over forward simulation because the apertures traverse non-physical sub-coordinate territory inaccessible to step-by-step kinetic evolution.
\end{remark}


\section{Discussion}
\label{sec:discussion}

\subsection{S-Entropy as Phase Space}

The results of this paper establish a precise analogy: \textbf{S-entropy is to categorical mechanics what phase space is to Newtonian mechanics}. In Newtonian mechanics, the laws of motion are expressed as differential equations on phase space $(q, p) \in T^*Q$~\cite{arnold1978mathematical}. The positions and momenta are not the physics---they are the coordinate system in which the physics becomes expressible as smooth differential equations. Without phase space, Newton's laws are algebraic relations on discrete events (collision outcomes, orbital periods); with phase space, they become a continuous flow amenable to the full apparatus of differential geometry.

Similarly, the discrete categorical structure ($\mathcal{O} = \mathcal{C} = \mathcal{P}$ at each node of a ternary hierarchy) is the physics. The S-entropy coordinates $(S_k, S_t, S_e)$ are not the physics---they are the coordinate system in which backward navigation becomes geodesic flow. Without S-entropy, backward navigation requires $\Omega(N)$ discrete enumeration. With S-entropy, it becomes $\mathcal{O}(\log_3 N)$ geodesic descent.

\subsection{Sufficient Statistics and Information Compression}

The S-entropy coordinates are sufficient statistics for the categorical state, in the precise sense of Fisher~\cite{fisher1922mathematical}: knowing $(S_k, S_t, S_e)$ provides the same probability of successful backward navigation as knowing the full discrete state. This sufficiency is not approximate---it is exact, following from the injectivity of the embedding (Theorem~\ref{thm:embedding_structure}) and the completeness of the geodesic navigation algorithm (Theorem~\ref{thm:geodesic_complexity}).

The compression from discrete state (a ternary string of length $k$, containing $k \log_2 3$ bits) to three continuous coordinates is lossless for navigational purposes. The continuous coordinates contain all the information needed for backward navigation, packaged in a form that admits gradient-based algorithms.

\subsection{Why $S = 0$ at the Solution}

A state with $S_k = 0$ is fully resolved: every level of the partition hierarchy has been determined. At this point, the categorical address is fully known, and no further navigation is possible or necessary. In this sense, $S_k = 0$ means ``nothing else can be known''---not because information is absent, but because the question has been completely answered. The solution forecloses all alternatives.

This connects to the Poincar\'{e} recurrence framework~\cite{poincare1890probleme}: the final state is the unique fixed point of the backward trajectory. The trajectory converges to $S_k = 0$ as the penultimate state is reached and the completion morphism is applied. The vanishing of $S_k$ marks the transition from navigation (ongoing, with uncertainty) to completion (finished, with certainty).

\subsection{Relationship to Optimal Transport}

The S-entropy geodesic has a natural interpretation in terms of optimal transport~\cite{villani2009optimal,kantorovich1942translocation}. The backward trajectory from $\mathbf{S}_f$ to $\mathbf{S}_0$ transports ``navigational mass'' from the final state to the root, and the geodesic is the path that minimizes the transport cost under the S-entropy metric. The Benamou--Brenier formulation~\cite{benamou2000computational} of optimal transport as geodesic flow in the Wasserstein space is structurally identical to our formulation of backward navigation as geodesic flow in S-entropy space.

This connection suggests that the S-entropy metric may be a specialization of the Wasserstein-2 metric to the ternary partition hierarchy---a direction for future investigation.

\subsection{Limitations}

Three limitations warrant acknowledgment:

\textbf{Ergodicity assumption.} The completeness of backward navigation relies on the Poincar\'{e} recurrence theorem, which requires measure-preserving flow in bounded phase space. Systems with KAM tori or other invariant structures may have regions inaccessible to geodesic navigation.

\textbf{Finite precision.} The continuous embedding is exact in the mathematical sense but must be discretized for computational implementation. The discretization error is bounded by $3^{-k}$ at depth $k$, which is exponentially small but non-zero.

\textbf{Metric singularities.} The Fisher-type metric $g_{ii} = 1/(S_i(1-S_i))$ diverges at $S_i = 0$ and $S_i = 1$. Geodesics approaching these boundaries have infinite arc length, reflecting the physical reality that complete resolution ($S_k = 0$) or complete uncertainty ($S_k = 1$) are limit points approached asymptotically but never exactly attained. In practice, navigation terminates at finite depth $k$ with $S_k = 1/k > 0$.


\section{Conclusion}
\label{sec:conclusion}

We have proved that S-entropy coordinates $(S_k, S_t, S_e) \in [0,1]^3$ constitute the unique continuous embedding of discrete categorical structure that enables backward trajectory completion in $\mathcal{O}(\log_3 N)$ time.

The argument has three pillars:

\textbf{Necessity.} Without a continuous metric, backward navigation through a ternary hierarchy requires $\Omega(N)$ enumeration. The discrete metric provides no directional information, reducing backward search to forward brute force (Theorem~\ref{thm:discrete_lower}).

\textbf{Uniqueness.} Among all continuous embeddings preserving the refinement order, ternary symmetry, and triple equivalence, S-entropy is the unique solution up to isometry. The three coordinates are forced by the three independent descriptions (partition, oscillator, category), and the metric coefficients are forced by the Fisher information structure and entropy consistency (Theorem~\ref{thm:uniqueness}).

\textbf{Self-similarity.} The S-entropy metric is recursively isometric: each sub-coordinate inherits the identical metric structure, eliminating any distinguished hierarchical level. The framework is scale-free---the same navigation algorithm applies at every level of description, and the level cannot be determined from intrinsic metric data alone (Theorem~\ref{thm:recursive_metric}).

The central metaphor is precise: S-entropy is to categorical mechanics what phase space is to Newtonian mechanics. The discrete triple-equivalence structure ($\mathcal{O} = \mathcal{C} = \mathcal{P}$ at every node) is the physics. S-entropy coordinates are the space in which that physics becomes navigable. The pendulum with three states already contains the entire framework; the rest is generalization to $3^k$ states through the recursive self-similarity of the embedding.

\bibliographystyle{unsrt}
\bibliography{references}

\end{document}
